% !TeX program = pdflatex
\PassOptionsToPackage{final}{graphicx}

\documentclass[preprint,12pt]{elsarticle}

% ---------- math ----------
\usepackage{amsmath}
\usepackage{amssymb}
\usepackage{bm}

% ---------- figures ----------
\usepackage{graphicx}
\DeclareGraphicsExtensions{.pdf,.png,.jpg,.jpeg}
% Overleaf-safe figure search paths:
% - ./        : manuscript.tex and figures/ are uploaded at project root
% - paper/   : project root contains paper/manuscript.tex and paper/figures/
% - ../paper/: local/build-directory fallback
\graphicspath{{./}{paper/}{./paper/}{../}{../paper/}}
\usepackage[table]{xcolor}

% ---------- tables ----------
\usepackage{booktabs}
\usepackage{array}
\usepackage{tabularx}
\usepackage{siunitx}

% ---------- links ----------
\usepackage{hyperref}
% \usepackage[hidelinks]{hyperref}
\usepackage{makecell}

% ---------- review ----------
\IfFileExists{lineno.sty}{%
  \usepackage{lineno}
  \linenumbers
}{%
  \providecommand{\linenumbers}{}
}

\usepackage{textcomp}

\IfFileExists{placeins.sty}{%
  \usepackage{placeins}
}{%
  \providecommand{\FloatBarrier}{}
}

% Keep figures and tables within the heading block where they are introduced.
\let\origsection\section
\renewcommand{\section}{\FloatBarrier\origsection}
\let\origsubsection\subsection
\renewcommand{\subsection}{\FloatBarrier\origsubsection}
\let\origsubsubsection\subsubsection
\renewcommand{\subsubsection}{\FloatBarrier\origsubsubsection}

\journal{Reliability Engineering \& System Safety}

\biboptions{authoryear,round}

\expandafter\def\csname appreftext:app:sph_density\endcsname{Appendix~A}
\expandafter\def\csname appreftext:sm:hazard_sensitivity\endcsname{Appendix~B}
\expandafter\def\csname appreftext:sm:disaster_locations\endcsname{Appendix~C}
\expandafter\def\csname appreftext:sm:spatial_diagnostics\endcsname{Appendix~D}
\DeclareRobustCommand{\appref}[1]{%
  \hyperref[#1]{%
    \ifcsname appreftext:#1\endcsname%
      \csname appreftext:#1\endcsname%
    \else%
      Appendix~\ref*{#1}%
    \fi%
  }%
}

\newif\ifanonymous
\anonymousfalse

\immediate\write18{texcount -inc -sum manuscript.tex > wordcount.txt}

% \renewcommand{\thesection}{\Alph{section}}

\begin{document}

\begin{frontmatter}

\title{Coupled Density--Hazard Navigation for Evacuation Reliability in Dense Urban Informal Settlements}
% \title{Risk-Oriented Evacuation Reliability Analysis in Dense Urban Informal Settlements Using Density--Hazard Coupled Navigation Fields}

\ifanonymous
% \author{Anonymous Submission}
\else
\author[aff1]{Chun Song}
\ead{csong945@connect.hkust-gz.edu.cn}
\author[aff2,aff3]{Xuchuan Lin}
\ead{linxuchuan@iem.ac.cn}
\author[aff1]{Rui Cao\corref{cor1}}
\ead{ruicao@hkust-gz.edu.cn}
\cortext[cor1]{Corresponding author.}

\affiliation[aff1]{
  organization={Thrust of Urban Governance and Design, Society Hub, The Hong Kong University of Science and Technology (Guangzhou)},
  city={Guangzhou},
  country={China}
}

\affiliation[aff2]{
  organization={Key Laboratory of Earthquake Engineering and Engineering Vibration, Institute of Engineering Mechanics, China Earthquake Administration},
  city={Harbin},
  country={China}
}

\affiliation[aff3]{
  organization={Key Laboratory of Earthquake Disaster Mitigation, Ministry of Emergency Management},
  city={Harbin},
  country={China}
}

\fi

\begin{abstract}
Dense urban informal settlements pose a distinctive evacuation-reliability challenge: extreme building and population densities, maze-like alley networks, and limited route redundancy generate heavy-tailed exposure outcomes beyond the scope of conventional guidance. Traditional static or density-only navigation treats congestion and hazard as separate processes, limiting its ability to ensure reliable clearance or control high-intensity exposure. To address this gap, we develop a coupled density–hazard dynamic navigation framework designed to enhance evacuation reliability in such environments. The framework estimates crowd density using a Smoothed Particle Hydrodynamics (SPH) field with a calibrated speed–density relationship, models hazard evolution as an arrival-time intensity field, and dynamically updates a global navigation potential integrating travel impedance, hazard penalties, and feasibility with density-aware local motion. We further propose a staged, tail-sensitive validation approach to assess the reliability of the framework. In the Shipai Village case, DH-60 retains the high-quantile clearance improvements achieved by dynamic routing while reducing the two model-based exposure indicators more than Static and density-only routing. These findings underscore that explicit hazard awareness—rather than congestion redistribution alone—is the key mechanism for controlling simulated exposure. The source code is available at \url{https://github.com/rui-research/density-hazard-navigation}.
\end{abstract}


\begin{keyword}
Informal settlements \sep
Evacuation routing \sep
Density--hazard coupling \sep
Pedestrian dynamics \sep
Dynamic navigation potential \sep
Tail-risk reliability analysis
\end{keyword}
\end{frontmatter}


\section{Introduction}
\label{sec:intro}

Urban villages are widespread informal settlements in rapidly
urbanising Chinese cities, characterised by dense population, complex morphology, and substandard living conditions \citep{cao2025mapping}. Their irregular building
shapes, ``handshake-house'' alleys, limited visibility, and frequent
dead-ends create labyrinthine outdoor pedestrian networks with weak path
redundancy~\citep{Al2014Villages_UHPress,Lin2011GuangzhouVIC}. During an
emergency, local congestion, blockage, smoke, or other hazard propagation
can cascade into severe queueing, local reversals, and trapping. The relevant
safety question therefore concerns not only average clearance, but also the
time by which the slowest agents evacuate and the exposure they
accumulate---the right tails of the evacuation-time and model-risk
distributions.

Existing evacuation research offers a rich methodological toolbox.
Microscopic models---social-force \citep{Helbing1995SocialForce},
cellular-automata and floor-field formulations
\citep{BlueAdler2001CA,Burstedde2001FloorField}---reproduce local
interactions and bottleneck phenomena, while continuum formulations
\citep{Hughes2002Continuum} support large-population flow analysis. A
parallel line of work develops state-responsive guidance and risk-aware
assessment: real-time rerouting under disruptions
\citep{Darvishan2021DynamicReroute}, GIS-based accessibility analysis
\citep{Liu2022EvacAccessibility}, BIM/FDS-coupled indoor fire evacuation
\citep{Ding2024BIMFireEvac}, dynamic fire navigation that jointly accounts for
travel time, fire impacts, and crowd density
\citep{Golshani2025FireNavigation}, and risk-perceptive path planning under
toxic or multi-hazard exposure
\citep{He2024RiskPerceptionPath,Luan2025ToxicGasAirport}. These studies establish
that congestion and hazard information can be combined in responsive routing.
The remaining gap concerns deployment and evaluation at outdoor settlement
scale: transferring joint state-responsive guidance from buildings and
facility networks to an informal-settlement raster with tens of thousands of
agents and hundreds of irregular corridors, while evaluating high-quantile
clearance and exposure outcomes across repeated scenarios.

Joint modelling matters because the two channels can act in opposition.
Congestion-aware routing redistributes flow away from crowded corridors,
but when congestion and hazard zones spatially correlate---as they often
do in narrow informal alleys where smoke or debris concentrates exactly
where flow concentrates---redirection can steer pedestrians \emph{toward}
hazard, increasing exposure relative to a hazard-unaware baseline.
Hazard-aware routing reduces exposure but, if it ignores crowd density,
can manufacture secondary bottlenecks. Reliable evacuation guidance in
maze-like settlements therefore requires a formulation in which
density-induced travel impedance and hazard-induced risk jointly reshape
the global navigation potential throughout the simulation, rather than
either channel acting alone on a static or partial cost surface.

This study develops a density--hazard coupled dynamic navigation framework
on a high-resolution raster domain reconstructed from GIS data. The
framework continuously estimates spatiotemporal crowd density, maps
density to mobility cost via a calibrated speed--density relation,
represents hazard evolution through an intensity field, and recomputes a
global navigation potential at selected update intervals. To make the two
governing parameters---the hazard-sensitivity weight $\beta$ and the
field update period $\Delta t_u$---empirically defensible rather than
hand-tuned, we adopt a two-stage protocol: parameters are first calibrated
on a simplified network amenable to dense Monte Carlo sampling, then
deployed on the full Shipai Village benchmark in Guangzhou
($\sim$\SI{0.57}{km^{2}}, \num{40000} agents). This design supports a
controlled comparison of hazard-coupled and congestion-only guidance
under matched conditions, isolating the contribution of explicit hazard
awareness to tail-risk reduction in dense informal-settlement evacuation.

The contributions of this work are four-fold:
\begin{itemize}

    \item \textbf{A density--hazard coupled navigation framework.}
    We formulate a unified, online-updatable global cost surface that
    simultaneously integrates Smoothed Particle Hydrodynamics (SPH)-estimated crowd density with a dynamic
    hazard arrival-time field, resolving the congestion--hazard
    interdependence that static and density-only methods treat as
    independent. The resulting navigation potential adapts in real time
    to both crowd state and hazard evolution at settlement scale.

    \item \textbf{A tail-sensitive evaluation protocol for evacuation
    reliability.} We introduce an assessment methodology that supplements
    mean-clearance summaries with high-quantile clearance times, Soft Risk,
    and Hard Harm, each reported with bootstrap confidence intervals. The protocol is
    specifically designed for regimes where extreme delay and
    exposure---not average performance---govern safety outcomes.

    \item \textbf{A transferable two-stage calibration--deployment
    methodology.} We establish a systematic procedure that identifies the
    framework's governing parameters $(\beta,\,\Delta t_{u})$ on a
    tractable reduced network before transferring them to full-scale
    deployment, separating dense parameter screening from the computational
    cost of full-scale evaluation. The resulting operating point is then
    verified and, where necessary, can be recalibrated for the target geometry.

    \item \textbf{Mechanistic evidence that hazard awareness drives
    simulated exposure reduction.} Through full-scale
    simulations on the Shipai Village benchmark, we demonstrate that
    explicit hazard embedding in the navigation potential---not
    congestion redistribution alone---is the primary mechanism
    responsible for reducing Soft Risk and Hard Harm, and we
    quantify the clearance--exposure trade-off and diminishing-returns
    structure across competing strategy classes.

\end{itemize}

The remainder of the paper is organised as follows.
Section~\ref{sec:related} situates the work within the literature on
pedestrian dynamics, hazard-aware routing, and informal-settlement
evacuation.
Section~\ref{sec:method} develops the coupled density--hazard navigation
framework, and Section~\ref{sec:evaluation} introduces the tail-sensitive
evaluation protocol used throughout the study.
Section~\ref{sec:calibration_validation} calibrates the underlying
pedestrian model against empirical fundamental diagrams and validates it
on canonical benchmarks.
Building on this foundation, Section~\ref{sec:strategy_calibration}
identifies the framework's governing parameters on a simplified
    settlement network via parameter calibration and hazard-field response
    illustrations.
Section~\ref{sec:shipai} transfers the calibrated framework to the
full-scale Shipai Village benchmark and reports the main evacuation
results.
Section~\ref{sec:discussion} interprets the findings, examines
practical implications, and acknowledges limitations.
Section~\ref{sec:conclusion} summarises the principal contributions and
outlines directions for future research.




\section{Related Work}
\label{sec:related}

This section reviews three strands of literature relevant to the proposed
framework: pedestrian-dynamics models as global guidance mechanisms
(\S\ref{sec:rw_basic}), dynamic and risk-aware navigation
(\S\ref{sec:rw_nav}), and reliability-oriented evaluation together with
prior work on informal-settlement evacuation (\S\ref{sec:rw_eval}). For
each strand we identify what it does and does not provide for outdoor
settlement-scale evacuation under coupled congestion and hazard.
Section~\ref{sec:rw_position} then positions our framework against the
closest comparators.

\subsection{Pedestrian dynamics models as guidance fields}
\label{sec:rw_basic}

Microscopic, mesoscopic, and macroscopic pedestrian-dynamics models have
been refined for several decades. Social-force models reproduce local
interactions and emergent phenomena including bottleneck oscillations and
lane formation~\citep{Helbing1995SocialForce}; cellular-automata and
floor-field formulations gain efficiency and offer implicit global guidance
through transition rules~\citep{BlueAdler2001CA,Burstedde2001FloorField};
continuum formulations describe crowd flow at the macroscopic level and
support large-population simulation
\citep{Hughes2002Continuum,Treuille2006ContinuumCrowds}. More recent
learning-based work spans data-driven pedestrian-motion models,
rapid evacuation-risk prediction, and AI-assisted safety analysis
\citep{Berggold2026MLReview,Li2023ELMEvacRisk,Lu2026HumanAI}; these are
distinct applications and are not all adaptive navigation methods.

These models capture local movement well, but the route structure they
produce is, in most implementations, anchored to a guidance representation
that is either fixed (a precomputed shortest-path field, a static
floor-field exit attractor) or only weakly responsive to evolving network
conditions. In high-tortuosity outdoor settlements, where a single
saturated corridor can invalidate a previously attractive route, a stale
global cost surface no longer reflects the post-disruption network state,
leaving agents on routes that may no longer minimise travel cost or
exposure even when their local movement is faithfully resolved.

\subsection{Dynamic and risk-aware guidance: progress and remaining gaps}
\label{sec:rw_nav}

A separate line of work has begun to make global guidance state-responsive.
Dynamic potential fields incorporate density effects to support rerouting
under congestion~\citep{Kretz2009DynamicPotential}, and Eikonal-type
travel-time fields provide globally consistent gradients efficiently
solvable by fast-marching or fast-sweeping
methods~\citep{Sethian1996FastMarching,Zhao2005FastSweeping}. Hazard-aware
guidance has matured along several relevant axes. BIM/FDS-based systems use
dynamic fire information for in-building route updates
\citep{Ding2024BIMFireEvac}. Golshani and Fang integrate FDS, agent-based
simulation, a modified Dijkstra algorithm, and dynamic signage so that travel
time, fire impacts, and crowd density jointly affect recommended routes
\citep{Golshani2025FireNavigation}. Risk-perceptive path-planning studies use
dynamic toxic or multi-hazard risk to inform route and response decisions
\citep{He2024RiskPerceptionPath,He2024DynamicMultiHazardEvac,Luan2025ToxicGasAirport,Yang2025MetroStampede},
whereas related work evaluates safety-barrier performance or crowd injury
mechanisms without necessarily providing online route guidance
\citep{Yuan2022SafetyBarrierEvac,Guo2024StampedeTerrorist}.
Simulation--optimisation and reinforcement-learning approaches further
expand the design
space~\citep{Zhang2024SimOptGuidance,Huang2023DSA_DDQN}.

Two deployment limitations motivate our approach. First, methods that jointly
use hazard and congestion information have been demonstrated mainly in
buildings or facility-scale networks; their computational and behavioural
properties do not directly establish performance for an outdoor
settlement-scale regime with tens of thousands of agents distributed across
hundreds of irregular corridors. Second, outdoor and informal-settlement
studies generally use road-network rerouting, fixed shortest paths, or
scenario-level route restrictions rather than a repeatedly updated raster
cost-to-go field. The methodological contribution pursued here is therefore
not the first combination of congestion and hazard, but a scalable
settlement-scale formulation in which density-derived slowness, hazard
intensity, and blocking jointly update one global raster potential, coupled
with a tail-sensitive assessment protocol.

\subsection{Reliability-oriented evaluation and informal-settlement
benchmarks}
\label{sec:rw_eval}

System-safety analysis of evacuation should look beyond the mean. Tail
indicators, exceedance probabilities, and Conditional Value-at-Risk are
established tools for evaluating systems whose failures are dominated by
extreme outcomes~\citep{Rockafellar2002CVaR}. Evacuation studies address
uncertainty and adverse outcomes through radiological scenario modelling,
airport-flow optimisation, ship-layout analysis, sampling-convergence
studies, and disruption-aware network rerouting
\citep{Kim2023RadiologicalEvac,Ma2024AirportPassengerFlow,Wang2022RoRoEvac,Lovreglio2019SamplingConvergence,Darvishan2021DynamicReroute}.
These works motivate risk- and uncertainty-aware assessment, but they do
not all report high-quantile or CVaR measures. Explicit upper-tail
reporting therefore remains less common than mean clearance-time analysis,
particularly for dense settlement scenarios.

Tail-sensitive evaluation is most consequential in settings where
mean-based summaries are known to mask safety failures; dense informal
settlements are arguably such a setting, yet the empirical literature on
them is itself underdeveloped. Urban villages and analogous layouts have
been extensively characterised in urban
morphology~\citep{Al2014Villages_UHPress,Lin2011GuangzhouVIC,cao2025mapping},
but evacuation studies remain sparse, in part because community-scale
drills are infeasible and detailed digital reconstruction is
labour-intensive. Recent simulation studies have begun to address this:
\citet{Butters2025Kibera} simulate evacuation in Kibera using
agent-based methods on open spatial data, and
\citet{Kuligowski2025Humanitarian} examine humanitarian-corridor
evacuation in informal contexts. High-resolution mapping with
LiDAR~\citep{Liu2025LiDARInformal} and accessibility-based difficulty
estimation~\citep{Liu2022EvacAccessibility} are making richer
reconstructions feasible. These efforts establish that
informal-settlement evacuation can be simulated and analysed; what they
do not yet combine, on the same high-density irregular-network case, is
(a) global dynamic guidance, (b) coupled density--hazard cost surfaces,
and (c) a tail-sensitive evaluation protocol with bootstrap inference.

\subsection{Positioning of this work}
\label{sec:rw_position}

Table~\ref{tab:comparator} positions the framework against the closest
comparators along deployment scale, hazard--congestion treatment, and the
form of guidance. Existing indoor work already demonstrates joint dynamic
fire--congestion navigation~\citep{Golshani2025FireNavigation}, while other
building and facility studies use dynamic fire or multi-hazard information
for route and response optimisation
\citep{Ding2024BIMFireEvac,He2024DynamicMultiHazardEvac,Luan2025ToxicGasAirport}.
\citet{Kretz2009DynamicPotential} provides density-responsive pedestrian
guidance without hazard, and recent informal-settlement simulators model
large populations and, in the flood case, dynamic hazard interaction
\citep{Butters2025Kibera,Kuligowski2025Humanitarian}. The distinctive
combination evaluated in this paper is an outdoor settlement-scale raster
potential that is repeatedly updated from both density-derived slowness and
hazard state, together with high-quantile and exposure-tail comparisons under
matched repeated scenarios. The specific contributions are enumerated in
Section~\ref{sec:intro}.


\begin{table}[t]
\centering
\small
\caption{Positioning against the closest dynamic and risk-aware evacuation
works. Treatment describes how hazard and congestion enter the model;
guidance distinguishes online route updates from fixed or scenario-level
route choices.}
\label{tab:comparator}
\begingroup
\setlength{\tabcolsep}{4pt}
\renewcommand{\arraystretch}{1.12}
\begin{tabularx}{\linewidth}{@{}
>{\raggedright\arraybackslash}p{0.23\linewidth}
>{\raggedright\arraybackslash}p{0.19\linewidth}
>{\raggedright\arraybackslash}p{0.28\linewidth}
>{\raggedright\arraybackslash}X@{}}
\toprule
\textbf{Reference} & \textbf{Scale} & \textbf{Treatment} & \textbf{Guidance} \\
\midrule
Kretz~\citeyearpar{Kretz2009DynamicPotential}
    & Generic pedestrian scenario
    & density-responsive; no hazard
    & dynamic quickest-path potential \\
Ding et al.~\citeyearpar{Ding2024BIMFireEvac}
    & Subway station
    & dynamic fire and exit utilisation
    & real-time route optimisation \\
Golshani--Fang~\citeyearpar{Golshani2025FireNavigation}
    & Multiple building cases
    & joint fire, travel time, and crowd density
    & dynamic Dijkstra/signage updates \\
He et al.~\citeyearpar{He2024DynamicMultiHazardEvac}
    & Chemical-industry area
    & dynamic multi-hazard risk
    & risk-informed route/start-time strategy \\
Luan et al.~\citeyearpar{Luan2025ToxicGasAirport}
    & Airport
    & toxic-gas exposure and response-mode risk
    & risk assessment and response selection \\
Butters--Dawson~\citeyearpar{Butters2025Kibera}
    & Informal settlement
    & dynamic flood interaction and congestion
    & fixed shortest path to shelter \\
Kuligowski et al.~\citeyearpar{Kuligowski2025Humanitarian}
    & Humanitarian settlement
    & density and route-restriction scenarios; no propagated hazard
    & fixed/scenario-level route choice \\
\midrule
\textbf{This work} &
\textbf{Settlement-scale outdoor} &
\textbf{joint density-derived slowness, hazard, and blocking} &
\textbf{periodically updated raster potential} \\
\bottomrule
\end{tabularx}
\endgroup
\end{table}


\section{Methodology}
\label{sec:method}

We formulate evacuation in dense informal-settlement networks as an
online global navigation problem in which route choice adapts periodically
to evolving congestion and hazard conditions. The objective is not simply to reach
an exit, but to do so over a time-varying cost landscape that
simultaneously reflects mobility loss and safety deterioration.

The framework comprises six components, all defined on a unified raster
domain (Figure~\ref{fig:framework}):
(i)~a multi-channel raster representation of the environment
(\S\ref{sec:raster});
(ii)~a continuous SPH density field with a calibrated speed--density
relation (\S\ref{sec:density});
(iii)~an arrival-time-driven hazard field
(\S\ref{sec:hazard_field});
(iv)~a density--hazard coupled global navigation potential
(\S\ref{sec:potential});
(v)~a density-aware local interaction model
(\S\ref{sec:interaction_model});
and (vi)~a global--local blending rule for agent motion
(\S\ref{sec:motion_update}).
Section~\ref{sec:implementation} reports parameter values, and
Section~\ref{sec:strategy_calibration} examines sensitivity to the
guidance-weight configuration. Supporting checks for the SPH density
estimator and hazard-parameter model are provided in
\appref{app:sph_density} and \appref{sm:hazard_sensitivity}, respectively.

Table~\ref{tab:notation} summarises the principal notation. All spatial
fields are defined on a unified raster grid $\mathcal{G}$ with cell size
$r \times r$; $\mathbf{p}$ denotes a cell centre and $\mathbf{x}_i$ the
continuous position of agent~$i$.



\begin{table}[htbp]
\centering
\caption{Principal notation used in the methodology.}
\label{tab:notation}
\footnotesize
\begingroup
\setlength{\tabcolsep}{4pt}
\renewcommand{\arraystretch}{1.12}
\begin{tabularx}{\linewidth}{@{}
>{\raggedright\arraybackslash}p{0.34\linewidth}
>{\raggedright\arraybackslash}X@{}}
\toprule
\textbf{Symbol} & \textbf{Description} \\
\midrule
$\mathcal{G}$, $r$ & Raster grid; cell resolution (m) \\
\makecell[l]{$\Omega_{\mathrm{walk}},\Omega_{\mathrm{obs}}$\\
$\Omega_{\mathrm{exit}}$} & Walkable, obstacle, and exit masks \\
$\mathbf{p}$, $\mathbf{x}_i$ & Cell centre on the raster; continuous position of agent $i$ \\
$\rho_i$, $\rho(\mathbf{p},t)$ & Agent-level and grid-level crowd density \\
$h$, $\Delta_\rho$ & SPH smoothing length; auxiliary density-aggregation cell size \\
$v(\mathbf{p},t)$, $c_p(\mathbf{p},t)$ & Speed field and pedestrian-density-dependent slowness (mobility cost) \\
$\varepsilon_v$ & Slowness regulariser (m/s) \\
$\varepsilon_\Phi$, $\varepsilon_d$ & Potential-gradient and direction-normalisation regularisers \\
$v_f$, $\rho_0$, $n$ & Fundamental-diagram parameters \\
$T_h(\mathbf{p})$, $v_h(\mathbf{p})$ & Hazard arrival-time field; hazard propagation speed \\
$R(\mathbf{p},t)$, $B(\mathbf{p},t)$ & Continuous hazard intensity/cost channel and hard-blocking mask \\
\makecell[l]{$R_{\max}$, $\alpha_{\mathrm{pre}}$, $\tau_{\mathrm{pre}}$\\
$\tau_{\mathrm{rise}}$, $R_{\mathrm{block}}$} & Hazard saturation, pre-arrival leakage, growth, and blocking parameters \\
$\Phi(\mathbf{p},t)$, $\mathbf{g}(\mathbf{p},t)$ & Global navigation potential and direction \\
$\alpha$; $\beta$; $\gamma$ & Dimensionless efficiency weight; safety conversion coefficient (s/m); hard-blocking switch ($0$ or $+\infty$) \\
$\Delta t_u$ & Period of global potential updates \\
\makecell[l]{$\rho_{\mathrm{low}},\rho_{\mathrm{high}}$\\
$\rho_{\mathrm{ref}}$} & Density thresholds for blending and pressure \\
$\lambda_i(t)$ & Global--local blending coefficient at agent $i$ \\
\bottomrule
\end{tabularx}
\endgroup
\end{table}



\begin{figure}[t]
\centering
\includegraphics[width=1.0\linewidth]{figures/03_methodology/framework.pdf}
\caption{Framework of the proposed dynamic navigation model. Pedestrian
positions are projected onto the raster to construct the SPH density
field $\rho(\mathbf{p},t)$, which is mapped through the calibrated
fundamental diagram to a speed field $v(\mathbf{p},t)$ and a slowness
field $c_p(\mathbf{p},t)$. The subscript $p$ denotes the pedestrian
mobility-cost channel. In parallel, an arrival-time-driven hazard
model produces a continuous hazard field $R(\mathbf{p},t)$ and a
hard-blocking mask $B(\mathbf{p},t)$. These routing components are distinct
from the outcome metrics Soft Risk $\Sigma_s$ and Hard Harm $H$ defined in
\S\ref{sec:eval_exposure}. They are combined into a
single cost field that defines the global navigation potential
$\Phi(\mathbf{p},t)$. Agent motion is determined by a density-dependent
blend of the global direction $-\nabla\Phi$ and local interaction
forces; agent positions feed back into the density field, closing the
loop.}
\label{fig:framework}
\end{figure}

% =========================================================
\subsection{Raster representation}
\label{sec:raster}

Dense informal-settlement networks exhibit irregular connectivity,
frequent dead-ends, and narrow alleys. To preserve fine-grained outdoor
geometry and to enable robust neighbourhood-based computation, we
represent the environment as a high-resolution multi-channel raster with
spatial resolution $r$ (e.g.\ $0.1$\,m). All GIS vector layers are
rasterised into five channels:
\begin{enumerate}
    \item \textbf{Walkable mask} $\Omega_{\mathrm{walk}}$: outdoor corridors and open spaces.
    \item \textbf{Obstacle/building mask} $\Omega_{\mathrm{obs}}$: non-walkable regions.
    \item \textbf{Exit mask} $\Omega_{\mathrm{exit}} \subset \Omega_{\mathrm{walk}}$: destination and safe areas.
    \item \textbf{Hazard intensity field} $R(\mathbf{p},t)$ (\S\ref{sec:hazard_field}).
    \item \textbf{Blocking field} $B(\mathbf{p},t)$ (\S\ref{sec:hazard_field}).
\end{enumerate}
This unified raster supports consistent 8-neighbourhood operations,
GPU-friendly parallel updates, and online fusion of congestion and
hazard effects within a single navigation solver.

% =========================================================
\subsection{Density field and mobility cost}
\label{sec:density}

\subsubsection{SPH density estimation}
\label{sec:density_estimation}

To represent local crowd concentration under irregular and anisotropic
pedestrian configurations, we adopt a Smoothed Particle Hydrodynamics
(SPH) density estimator~\citep{Monaghan1992SPH,muller2003particle}.
Compared with cell counting on a fixed raster, SPH provides continuous
agent-level density estimates that remain robust in narrow alleys and
heterogeneous local layouts.

For agent $i$ at $\mathbf{x}_i$, the local density is
\begin{equation}
\rho_i = \sum_{j} m_j \, W_{\mathrm{poly6}}(\|\mathbf{x}_i-\mathbf{x}_j\|, h),
\label{eq:sph_density_agent}
\end{equation}
where $m_j$ is the particle weight ($m_j = 1$ in this study), $h$ is the
smoothing length, and $W_{\mathrm{poly6}}$ is the two-dimensional Poly6
kernel~\citep{muller2003particle}:
\begin{equation}
W_{\mathrm{poly6}}(r, h) = \frac{4}{\pi h^8}
\begin{cases}
(h^2-r^2)^3, & 0 \le r \le h,\\
0, & r > h.
\end{cases}
\label{eq:poly6_kernel}
\end{equation}
The kernel is compactly supported and satisfies
$\int_{\mathbb{R}^2} W_{\mathrm{poly6}}(r, h)\,\mathrm{d}A = 1$, so the
SPH estimate converges to the physical density under uniform sampling
away from boundaries.

For dynamic routing, the same agent-level estimates $\rho_i$ are mapped
to an auxiliary raster with cell width $\Delta_\rho$. Let
$\mathcal{A}_k(t)$ denote the active agents located in auxiliary cell
$k$. The cell density is
\begin{equation}
\rho_k^{\max}(t)=
\begin{cases}
\displaystyle\max_{i\in\mathcal{A}_k(t)}\rho_i(t),
& \mathcal{A}_k(t)\ne\varnothing,\\[4pt]
0, & \mathcal{A}_k(t)=\varnothing,
\end{cases}
\label{eq:sph_density_grid}
\end{equation}
and each fine navigation cell $\mathbf{p}$ inherits the value of its
containing auxiliary cell,
$\rho(\mathbf{p},t)=\rho_{k(\mathbf{p})}^{\max}(t)$. The implementation
uses $\Delta_\rho=1.0$\,m and nearest-cell upsampling to the $0.1$\,m
navigation raster. Thus there is no second Poly6 projection length
in the model: both agent motion and navigation originate from the single
Poly6 estimate with $h=1.2$\,m. The max aggregation conservatively
preserves the strongest local congestion signal within each auxiliary
cell.

\subsubsection{Sanity check of the SPH estimator}
\label{sec:density_check}

We verify that the continuous Poly6 estimator used to obtain $\rho_i$,
with the implemented support radius $h=1.2$\,m, recovers physical
pedestrian density under both uniform and random hard-core
configurations. Multi-centre averaging is applied to reduce spatial
variance. Across the tested density range ($0.1$--$4.5$\,ped/m$^2$),
the estimate matches the physical density to
\begin{equation}
\rho_{\mathrm{real}}
= 1.020\, \hat{\rho}_{\mathrm{SPH}} - 0.019,
\qquad R^2 = 1.000.
\end{equation}
This is a sanity check rather than a calibration in the modelling
sense---no free parameter of the framework is fitted here---and confirms
that the $h=1.2$\,m Poly6 estimate is suitable as the physical input to
the speed--density relation that follows. It validates the density
estimator, not the separate spatial fidelity of the deterministic
1.0\,m max-aggregation step in Eq.~\eqref{eq:sph_density_grid}. The full
protocol is reported in \appref{app:sph_density}.

\subsubsection{Speed--density relation and slowness field}
\label{sec:speed_density_cost}

Given $\rho(\mathbf{p}, t)$, we obtain a density-dependent desired speed
through a fundamental
diagram~\citep{weidmann1993fundamental,Helbing1995SocialForce} calibrated against
empirical pedestrian data (\S\ref{sec:calibration_fd}):
\begin{equation}
v(\mathbf{p}, t) = v_f \exp\!\left[ -\!\left( \frac{\rho(\mathbf{p}, t)}{\rho_0} \right)^{n} \right],
\label{eq:speed_field}
\end{equation}
with calibrated values $v_f = 1.574$\,m/s, $\rho_0 = 1.735$\,ped/m$^2$,
and $n = 0.852$. The corresponding slowness field
\begin{equation}
c_p(\mathbf{p}, t) = \frac{1}{\max\!\big(v(\mathbf{p}, t), \varepsilon_v\big)}
\label{eq:slowness_field}
\end{equation}
gives local mobility cost per unit traversed distance, with
$\varepsilon_v > 0$ avoiding singularity. The single $h=1.2$\,m SPH
estimate thus serves as a common physical basis for three coupled
modules: the fundamental diagram (speed adaptation), the aggregated
raster slowness field (navigation cost), and the local interaction force
(\S\ref{sec:interaction_model}).

% =========================================================
\subsection{Hazard field}
\label{sec:hazard_field}

We model hazard propagation as a front emanating from a set of source
cells $\mathcal{S}$ through a medium with spatially varying speed
$v_h(\mathbf{p})$, which depends on the underlying material (e.g.\ open
space versus building footprint). The hazard arrival time is the minimum
travel time from sources, satisfying the Eikonal
equation~\citep{Sethian1996FastMarching}:
\begin{equation}
T_h(\mathbf{p}) = \min_{\pi_h:\, \mathcal{S} \to \mathbf{p}}
\int_{\pi_h} \frac{1}{v_h(\boldsymbol{\xi})}\,\mathrm{d}\ell,
\label{eq:hazard_arrival}
\end{equation}
solved by fast marching on the same raster
$\mathcal{G}$~\citep{Sethian1996FastMarching,Zhao2005FastSweeping}.
$T_h$ is computed once per simulation provided that $v_h$ and the source
set are fixed; if either changes, the field is recomputed.

The arrival-time field sets the local clock for hazard exposure but does
not imply a discontinuous jump from zero intensity. Let
$s=t-T_h(\mathbf{p})$ denote time relative to front arrival and
$R_0=\alpha_{\mathrm{pre}}R_{\max}$ denote the arrival-level intensity.
We use a pre-arrival leakage branch followed by monotone post-arrival
growth:
\begin{equation}
R(\mathbf{p}, t) =
\begin{cases}
R_0 \exp\!\left(\dfrac{s}{\tau_{\mathrm{pre}}}\right), & s < 0,\\[5pt]
R_0 + (R_{\max}-R_0)
\left[1-\exp\!\left(-\dfrac{s}{\tau_{\mathrm{rise}}}\right)\right],
& s \ge 0.
\end{cases}
\label{eq:hazard_intensity}
\end{equation}
The hard-blocking mask is obtained by thresholding intensity at a
critical level $R_{\mathrm{block}}$ after front arrival:
\begin{equation}
B(\mathbf{p}, t)
= \mathbb{1}\!\big[\, s \ge 0 \;\wedge\;
R(\mathbf{p}, t) \ge R_{\mathrm{block}} \,\big].
\label{eq:hazard_blocking}
\end{equation}
$B$ is therefore a binary closure mask that activates only after front
arrival and is recomputed whenever $R$ crosses $R_{\mathrm{block}}$.
All calibration and Shipai experiments reported in this paper use the
dimensionless values $R_{\max}=10$ and $R_{\mathrm{block}}=6.0$.
$R$ is therefore a dimensionless model risk score, not a concentration,
heat flux, fractional effective dose, or dose--response probability.
$R_{\mathrm{block}}$ is the navigation-closure threshold and is distinct
from the lower outcome-evaluation threshold $R_H=4.5$ in
Eq.~\eqref{eq:threshold_exceedance}. Parameter values for $v_h$,
$\alpha_{\mathrm{pre}}$, $\tau_{\mathrm{pre}}$, and
$\tau_{\mathrm{rise}}$, together with one-factor illustrations of how
these parameters change the constructed hazard field, are reported in
\appref{sm:hazard_sensitivity}.

% =========================================================
\subsection{Density--hazard coupled global potential}
\label{sec:potential}

\subsubsection{Continuous formulation}
\label{sec:phi_continuous}

Let $\Phi(\mathbf{p}, t)$ denote the dynamic global potential at cell
$\mathbf{p}$ and time $t$, defined as the minimum cumulative arrival
cost to the nearest exit under time-varying mobility, hazard, and
blocking conditions:
\begin{equation}
\Phi(\mathbf{p}, t) = \min_{\pi:\, \mathbf{p} \to \Omega_{\mathrm{exit}}}
\int_\pi \Big[
\alpha\, c_p(\mathbf{s}, t)
+ \beta\, R(\mathbf{s}, t)
+ \gamma\, B(\mathbf{s}, t)
\Big]\,\mathrm{d}\mathbf{s},
\label{eq:phi_continuous}
\end{equation}
Here $R$ and $B$ are dimensionless. We set $\alpha$ as a dimensionless
weight and $\beta$ as a conversion coefficient with units of
\si{\second\per\meter} per unit of $R$. The blocking coefficient is not
a tunable finite penalty: $\gamma=0$ ignores $B$, whereas
$\gamma=+\infty$ imposes exact exclusion. In the latter case we use the
extended-real convention $\gamma B=0$ for $B=0$ and
$\gamma B=+\infty$ for $B=1$. Thus the integrand has units of
\si{\second\per\meter} on traversable cells, $\Phi$ has units of seconds,
and any path entering a blocked cell has infinite cost. Unlike
static distance-to-exit potentials, $\Phi$ jointly integrates
congestion-induced travel time, continuous hazard cost, and feasibility along
each candidate path, producing a single cost-to-go surface in which the
two channels co-determine route choice rather than acting as
independent penalties on a fixed network. Figure~\ref{fig:densityaware}
summarises how the density-driven slowness, continuous hazard cost, and hard
blocking terms are assembled into this dynamic cost field. The selection
and sensitivity of the weights are discussed in \S\ref{sec:implementation} and
\S\ref{sec:strategy_calibration}.



\begin{figure}[t]
\centering
\includegraphics[width=1.0\linewidth]{figures/03_methodology/potential.pdf}
\caption{Composition of the dynamic cost field. The density-driven
slowness $c_p$, continuous hazard cost $R$, and hard-blocking mask $B$ are
combined into a single cost surface whose Bellman shortest-arrival
solution gives the global potential $\Phi(\mathbf{p},t)$. The
navigation direction is the normalised negative gradient of $\Phi$,
allowing pedestrians to avoid both congestion and hazardous regions
while moving toward exits. Here $R$ and $B$ are routing inputs; Soft Risk
$\Sigma_s$ and Hard Harm $H$ are evaluation outcomes and are not additional
terms in the potential.}
\label{fig:densityaware}
\end{figure}


\subsubsection{Discrete update on the raster}
\label{sec:phi_discrete}

Equation~\eqref{eq:phi_continuous} is realised on the raster as a
shortest-arrival-cost transform~\citep{Sethian1996FastMarching}. For
neighbouring cells $\mathbf{p}$ and
$\mathbf{q} \in \mathcal{N}_8(\mathbf{p})$ with step length
$\Delta s_{\mathbf{pq}}$, the edge cost at time $t$ is
\begin{equation}
c(\mathbf{p} \to \mathbf{q};\, t) =
\begin{cases}
+\infty,
& \gamma=+\infty \ \wedge\ B(\mathbf{q},t)=1,\\[4pt]
\Delta s_{\mathbf{pq}}
\big[\alpha c_p(\mathbf{q},t)+\beta R(\mathbf{q},t)\big],
& \text{otherwise},
\end{cases}
\label{eq:edge_cost}
\end{equation}
which has units of seconds, consistent with the line integral in
\eqref{eq:phi_continuous}. For Density--Hazard guidance,
$\gamma=+\infty$ is implemented by removing cells with $B=1$ from the
walkable graph; they are strictly impassable, rather than merely
discouraged by a large finite cost. Strategies with $\gamma=0$ ignore
$B$ in navigation.
The discrete potential satisfies the Bellman optimality equation
\begin{equation}
\Phi(\mathbf{p}, t) =
\begin{cases}
0, & \mathbf{p} \in \Omega_{\mathrm{exit}},\\[4pt]
\displaystyle
\min_{\mathbf{q} \in \mathcal{N}_8(\mathbf{p})}
\Big\{ \Phi(\mathbf{q}, t) + c(\mathbf{p} \to \mathbf{q};\, t) \Big\},
& \mathbf{p} \in \Omega_{\mathrm{walk}},\\[8pt]
+\infty, & \mathbf{p} \notin \Omega_{\mathrm{walk}}.
\end{cases}
\label{eq:phi_discrete}
\end{equation}
The Moore (8-connected) neighbourhood $\mathcal{N}_8(\mathbf{p})$
better approximates isotropic propagation than a 4-neighbourhood and
improves stability in narrow, irregular corridors.
Equation~\eqref{eq:phi_discrete} is solved by a label-correcting sweep
until convergence; with non-negative edge weights and hard $B$, the
solve is equivalent to a Dijkstra computation on the raster graph.

The global navigation direction is the normalised negative gradient of
the potential,
\begin{equation}
\mathbf{g}(\mathbf{p}, t) = -\frac{\nabla\Phi(\mathbf{p}, t)}
{\|\nabla\Phi(\mathbf{p}, t)\| + \varepsilon_\Phi},
\label{eq:nav_dir}
\end{equation}
with $\nabla\Phi$ approximated by central finite differences on the
raster.

\subsubsection{Online update strategy}
\label{sec:online_update}

Recomputing $\Phi(\cdot, t)$ at every simulation step is unnecessary,
since $\rho$ and $R$ vary on time scales much longer than the
integration step $\Delta t$. The framework supports periodic updates at
a configurable interval $\Delta t_u$ and, in principle, event-triggered
updates when $\rho$ or $R$ changes beyond a prescribed threshold. All
reported experiments use the periodic mode, with $\Delta t_u$ calibrated
in \S\ref{sec:strategy_calibration}. Independently of the trigger mode,
the raster relaxation can be GPU-parallelised. The cost of a single potential update on the Shipai
benchmark is reported in \S\ref{sec:strategy_calibration}.



% =========================================================
\subsection{Density-aware local interaction}
\label{sec:interaction_model}

While the Poly6 kernel is used for density estimation, short-range
interaction forces are computed using the Spiky
kernel~\citep{muller2003particle}, which provides a steeper and
numerically robust gradient for repulsive force evaluation:
\begin{equation}
W_{\mathrm{spiky}}(r, h) = \frac{10}{\pi h^5}
\begin{cases}
(h - r)^3, & 0 \le r \le h,\\
0, & r > h,
\end{cases}
\label{eq:spiky_kernel}
\end{equation}
with gradient (for $0 < r_{ij} \le h$)
\begin{equation}
\nabla_i W_{\mathrm{spiky}}(r_{ij}, h)
= -\frac{30}{\pi h^5}\,(h - r_{ij})^2\,
\frac{\mathbf{x}_i - \mathbf{x}_j}{r_{ij}},
\label{eq:spiky_gradient}
\end{equation}
where $r_{ij} = \|\mathbf{x}_i - \mathbf{x}_j\|$.

The SPH-estimated density is mapped to a pressure-like quantity that
modulates interaction strength:
\begin{equation}
P_i = k_p\, (\rho_i - \rho_{\mathrm{ref}})_+,
\qquad (\cdot)_+ = \max(\cdot, 0),
\label{eq:pressure_mapping}
\end{equation}
with activation threshold $\rho_{\mathrm{ref}}$ below which repulsion is
inactive, and stiffness $k_p$. The interaction force on pedestrian $i$
is
\begin{equation}
\mathbf{f}_i^{\mathrm{int}}
= -\sum_j m_j\!
\left( \frac{P_i}{\rho_i^2} + \frac{P_j}{\rho_j^2} \right)
\nabla_i W_{\mathrm{spiky}}(r_{ij}, h).
\label{eq:interaction_force_spiky}
\end{equation}
This formulation produces strong short-range repulsion in high-density
zones, capturing congestion build-up, bottleneck pressure, and
stop-and-go effects, while avoiding the particle clustering that can
arise when the density kernel itself is reused for force computation.

Compared with classical social-force
models~\citep{Helbing1995SocialForce}, in which repulsive forces depend
solely on pairwise distance, the proposed formulation derives density,
pressure, and force within a unified SPH framework using the Poly6 density
kernel and Spiky interaction kernel. The resulting agent-level density
estimates are quantitatively comparable with empirical density, allowing consistent integration with
raster-based routing and the calibrated speed--density relation
in~\eqref{eq:speed_field}.

% =========================================================
\subsection{Global--local blending and motion update}
\label{sec:motion_update}

Agent motion combines the global navigation cue with local
interaction-driven adaptation. Let
\begin{equation}
\mathbf{g}_i(t) = \mathbf{g}(\mathbf{x}_i, t)
\label{eq:global_dir}
\end{equation}
be the global direction at the agent's position, and let
$\mathbf{l}_i(t) = f(\mathcal{N}_i)$ be a local direction derived from
interaction forces (collision avoidance, lateral adjustment, alignment
in narrow corridors). The density-dependent blending coefficient
\begin{equation}
\lambda_i(t) = \mathrm{clip}\!\left(
\frac{\rho_i(t) - \rho_{\mathrm{low}}}
     {\rho_{\mathrm{high}} - \rho_{\mathrm{low}}};\;
0,\; 1\right)
\label{eq:lambda_blend}
\end{equation}
increases the weight of local interaction under congestion, and the two
cues are combined as
\begin{equation}
\mathbf{d}_i(t) = (1 - \lambda_i(t))\, \mathbf{g}_i(t)
+ \lambda_i(t)\, \mathbf{l}_i(t).
\label{eq:blend_dir}
\end{equation}
In free flow ($\rho_i < \rho_{\mathrm{low}}$) the agent follows global
guidance; in saturation ($\rho_i > \rho_{\mathrm{high}}$) local forces
dominate, modelling that an agent compressed against neighbours cannot
follow guidance against the resulting contact pressure. This trade-off
can in principle leave isolated agents trapped inside dense clusters; we
examine its consequence on tail evacuation behaviour empirically in
\S\ref{sec:shipai}.

The desired speed magnitude follows the same fundamental diagram used
for the slowness field,
\begin{equation}
v_i(t) = v\!\big(\rho_i(t)\big),
\label{eq:desired_speed}
\end{equation}
and velocity and position updates use forward Euler integration with
step $\Delta t$:
\begin{equation}
\mathbf{v}_i(t) = v_i(t)\,
\frac{\mathbf{d}_i(t)}{\|\mathbf{d}_i(t)\| + \varepsilon_d},
\qquad
\mathbf{x}_i(t + \Delta t) = \mathbf{x}_i(t) + \mathbf{v}_i(t)\, \Delta t.
\label{eq:state_update}
\end{equation}
This global--local coupling yields globally coherent routing in
maze-like corridors while capturing local density-driven adaptation at
bottlenecks, and naturally supports online hazard-aware rerouting
through the shared dynamic potential field.

% =========================================================
\subsection{Implementation}
\label{sec:implementation}

This section reports the parameter values used in the proposed
framework. Unless otherwise stated, all parameters are fixed across all
simulations; the two parameters subjected to formal calibration on a
simplified network ($\beta$ and $\Delta t_u$) are reported in
\S\ref{sec:strategy_calibration}.

The Poly6 smoothing length is $h=1.2$\,m in the density estimator,
neighbour search, and local interaction calculation. The resulting
agent densities are max-aggregated on an auxiliary raster with
$\Delta_\rho=1.0$\,m and upsampled by nearest-cell assignment to the
$r=0.1$\,m navigation raster; no second kernel length is used. The
simulation step is $\Delta t = 0.1$\,s, sufficient for
stable forward Euler integration; a small constant
$\varepsilon_v = \SI{1e-3}{\meter\per\second}$ avoids division by zero in
low-density regions.
Separate positive regularisers $\varepsilon_\Phi$ and $\varepsilon_d$ are
used in the potential-gradient and direction normalisations, respectively,
so that quantities with different dimensions do not share one symbol.
Density thresholds for the blending and pressure laws are
$\rho_{\mathrm{low}} = 0.5$\,ped/m$^2$,
$\rho_{\mathrm{high}} = 3.0$\,ped/m$^2$, and
$\rho_{\mathrm{ref}} = 1.0$\,ped/m$^2$, with stiffness $k_p = 500$. The
fundamental-diagram parameters $(v_f, \rho_0, n)$ are calibrated in
\S\ref{sec:calibration_fd}. The efficiency weight $\alpha$ is normalised
to $1$ throughout. The other guidance parameters are strategy-specific:
Static and Density-aware routing both use $\beta=\gamma=0$, whereas the
Density--Hazard strategy calibrates $\beta$ in
\si{\second\per\meter} and uses $\gamma=+\infty$ to remove cells with
$B=1$ from the walkable graph.
The update period $\Delta t_u$ is
calibrated in \S\ref{sec:strategy_calibration}. Hazard parameters
$(v_h, R_{\max}, \alpha_{\mathrm{pre}}, \tau_{\mathrm{pre}},
\tau_{\mathrm{rise}}, R_{\mathrm{block}})$ are reported in
\appref{sm:hazard_sensitivity}. The hazard field is simulated for all
strategies so that exposure outcomes can be evaluated consistently, but
neither $R$ nor $B$ enters the navigation calculation for Static or
Density-aware routing.

The normalisation $\alpha=1$ is a scale choice rather than a fitted
parameter. The coefficient $\beta$ converts one unit of the model-defined
hazard score into an equivalent time cost per metre. Thus
$\beta=\SI{1}{\second\per\meter}$ means that increasing $R$ by one unit
adds \SI{1}{\second\per\meter} to the local path-cost density; traversing
a length $L$ at a uniform intensity $R$ adds $R L$ seconds to the path
cost. This is a routing trade-off, not a claim that one hazard unit equals
one second of physical exposure or toxicity. Under the calibrated
speed--density relation,
$c_p(0)=1/v_f=0.64$\,s/m, $c_p(3.0)=3.13$\,s/m at the saturation
threshold $\rho_{\mathrm{high}}$, and severe crowding at
$\rho\approx 5.7$\,ped/m$^2$ gives $c_p\approx 10$\,s/m. Since the
hazard intensity is bounded by $R_{\max}=10$, the density-induced
slowness term with $\alpha=1$ and the saturated hazard term with
$\beta=\SI{1}{\second\per\meter}$ both contribute approximately
\SI{10}{\second\per\meter} in the high-density, saturated-hazard limit.
This comparison provides the scale rationale for the tested $\beta$
range; the selected value remains an empirical hazard-aversion
coefficient and must be recalibrated if the definition or numerical scale
of $R$ changes. Hard blocking is not approximated by a finite conversion
coefficient: $\gamma=+\infty$ assigns infinite edge cost to $B=1$ cells,
which the implementation removes from the walkable graph.
Table~\ref{tab:params_summary} summarises all values.

\begin{table}[htbp]
\centering
\caption{Model and strategy parameters. Calibrated and
strategy-specific values are reported in the indicated section.}
\label{tab:params_summary}
\small
\begin{tabularx}{\linewidth}{@{}>{\raggedright\arraybackslash}p{0.30\linewidth}
>{\centering\arraybackslash}p{0.25\linewidth}
>{\raggedright\arraybackslash}X@{}}
\toprule
\textbf{Parameter} & \textbf{Value} & \textbf{Description} \\
\midrule
Smoothing length $h$            & $1.2$\,m   & Poly6 density and interaction support \\
Density aggregation $\Delta_\rho$ & $1.0$\,m & Auxiliary max-aggregation raster \\
Navigation raster $r$           & $0.1$\,m   & Fine potential-field resolution \\
Time step $\Delta t$            & $0.1$\,s   & Forward Euler step \\
$\rho_{\mathrm{low}}$           & $0.5$\,ped/m$^2$ & Free-flow threshold \\
$\rho_{\mathrm{high}}$          & $3.0$\,ped/m$^2$ & Saturation threshold \\
$\rho_{\mathrm{ref}}$           & $1.0$\,ped/m$^2$ & Pressure activation \\
Pressure stiffness $k_p$        & $500$      & Repulsion strength \\
Slowness regulariser $\varepsilon_v$ & \SI{1e-3}{\meter\per\second} & Lower bound in $c_p=1/\max(v,\varepsilon_v)$ \\
$\varepsilon_\Phi$, $\varepsilon_d$ & Positive & Dimensionally separate normalisation regularisers \\
Free-flow speed $v_f$           & $1.574$\,m/s    & FD calibration, \S\ref{sec:calibration_fd} \\
$\rho_0$                        & $1.735$\,ped/m$^2$ & FD calibration, \S\ref{sec:calibration_fd} \\
$n$                             & $0.852$    & FD calibration, \S\ref{sec:calibration_fd} \\
$\alpha$                        & $1$        & Dimensionless efficiency weight (normalised) \\
$\beta$                         & $0$ or calibrated (s/m) & Hazard-score-to-time-cost conversion \\
$\gamma$                        & $0$ or $+\infty$ & Ignore $B$ or impose exact hard blocking \\
$R_{\mathrm{block}}$           & $6.0$      & Dimensionless navigation-closure threshold \\
$\Delta t_u$                    & $\infty$ or calibrated & Static or dynamic-field update period, \S\ref{sec:strategy_calibration} \\
\bottomrule
\end{tabularx}
\end{table}

% =========================================================
% Section 4: Evaluation Metrics and Assessment Protocol
% Detailed spatial diagnostics are provided in Appendix D.
% =========================================================
\FloatBarrier
\section{Evaluation Metrics and Assessment Protocol}
\label{sec:evaluation}

Evaluation is designed to assess evacuation performance and model-based
exposure outcomes. The reported outcome metrics are organised into three
families---efficiency (\S\ref{sec:eval_efficiency}), Soft Risk and Hard Harm
(\S\ref{sec:eval_exposure}), and congestion
(\S\ref{sec:eval_congestion})---together with protocol-level reliability
and tail-risk definitions (\S\ref{sec:eval_reliability}). They are applied
through a unified assessment protocol (\S\ref{sec:assessment_protocol}). Because later
sections also report local calibration, computation, and spatial
diagnostics, Table~\ref{tab:metric_summary} serves as a compact index of
the quantities used across the analysis rather than as a list of core
outcome metrics only.

\begin{table}[htbp]
\centering
\small
\setlength{\tabcolsep}{6pt}
\renewcommand{\arraystretch}{1.08}
\caption{Compact index of core metrics and auxiliary diagnostics used
across the calibration, strategy-evaluation, and Shipai analyses.
Spatial diagnostics ($\dagger$) are detailed in
\appref{sm:spatial_diagnostics}; local validation and computation
diagnostics are defined where they are used.}
\label{tab:metric_summary}
\begin{tabularx}{\textwidth}{@{}
  >{\raggedright\arraybackslash}p{0.44\textwidth}
  >{\raggedright\arraybackslash}X@{}}
\toprule
\textbf{Metric / diagnostic} & \textbf{Meaning} \\
\midrule
RMSE, $R^{2}$; $J_{\max}, C$; $C^{\mathrm{sim}}, C^{\mathrm{emp}}$; rel.\ error
  & Fit quality and flow-capacity validation. \\
$T_p^{(k)}$; $\mu_K[T_p]$, $Q_{q,K}[T_p]$
  & Within-run tail completion time and explicitly labelled cross-run summaries. \\
$E_i^{(k)}$; $\Sigma_s^{(k)}$
  & Agent exposure and population Soft Risk accumulations of the model risk score
    (risk-score--time-step and risk-score--person--time-step units). \\
$H^{(k)}$; $\mu_K[H]$
  & Run-level and cross-run-mean Hard Harm tallies
    (agent--time-step comparative threshold-event proxies). \\
$M^{+}_{8}, N^{+}_{8}$
  & Update-period congestion diagnostics. \\
Total computation; field-update time/fraction
  & Runtime cost of online field recomputation. \\
$\mathcal{R}_T(\tau),\ \mathcal{R}_E(\xi)$;
$\operatorname{VaR}_{0.95},\ \operatorname{CVaR}_{0.95}$
  & Protocol-level reliability and tail-risk quantities. \\
$F_{\mathrm{route}}(\mathbf{p}), L_{\rho}(\mathbf{p})$\,$^\dagger$;
$T_{\mathrm{cong}}(\mathbf{p}), L_{\mathrm{cong}}(\mathbf{p})$\,$^\dagger$
  & Spatial redistribution and congestion diagnostics. \\
Positive-cell quantiles; same-field Static-threshold counts$^\dagger$
  & Static-relative spatial-reduction summaries. \\
\bottomrule
\end{tabularx}
\end{table}

Spatial diagnostics used to localise persistent bottlenecks and
exposure hotspots are defined in \appref{sm:spatial_diagnostics}. In the
main text they are used only to interpret residual differences between
strategies; the primary comparisons rely on tail clearance, Soft Risk,
Hard Harm, and explicitly identified ensemble summaries.






% ---------------------------------------------------------
\subsection{Efficiency}
\label{sec:eval_efficiency}

Let $T_i$ denote the egress time of agent~$i$. We report completion
quantiles $T_{0.95}$ and $T_{0.99}$, the times by which $95\%$ and
$99\%$ of agents have evacuated, respectively. An agent still active at
the simulation horizon has a right-censored egress time and is not treated
as having exited; if the required completion fraction is not reached, the
corresponding completion quantile is reported as right-censored. High-quantile clearance is
more informative than the mean for safety analysis because mean
clearance can be dominated by a fast-moving majority that evacuates
well before the agents most at risk; in maze-like settlements, the
right tail of the egress-time distribution is precisely where
infrastructure design and guidance strategy decisions matter most.


% ---------------------------------------------------------
\subsection{Soft Risk and Hard Harm}
\label{sec:eval_exposure}

We retain the implementation-facing metric names \emph{Soft Risk} and
\emph{Hard Harm}, but use them only as labels for comparative model
outputs. Soft Risk $\Sigma_s$ accumulates the full modelled risk-score
history, whereas Hard Harm $H$ counts threshold-exceeding agent--time-step
samples. Here, ``hard'' denotes thresholding rather than observed physical
injury. Neither metric is a calibrated physical dose, affected-person
count, casualty count, or probability of harm.

Let $r_{i,n}=R(\mathbf{x}_i(t_n),t_n)$ be the dimensionless risk score,
measured in model risk-score units (RSU),
sampled at the nearest raster cell occupied by active agent $i$ at
$t_n=n\Delta t$. Let $\mathcal{A}_n$ denote the set of agents still active
at $t_n$. The quantity stored by the simulation for one agent is
\begin{equation}
E_i = \sum_{n:\,i\in\mathcal{A}_n} r_{i,n},
\label{eq:exposure}
\end{equation}
in risk-score--time-step units. Exposure for an unevacuated agent is
accumulated through the final simulated step without treating the censored
time as an exit. The run-level population Soft Risk accumulator is
\begin{equation}
\Sigma_s = \sum_i E_i
= \sum_n\sum_{i\in\mathcal{A}_n} r_{i,n},
\label{eq:soft_risk_population}
\end{equation}
with units of risk-score--person--time-step. This is the exact operation
implemented in the CUDA kernel: one local value of $R$ is added for each
active agent at every simulation step. Multiplying by $\Delta t$ gives
the time-integrated score $\Delta t\,\Sigma_s$ in
risk-score--person--s. A normalised full-intensity-equivalent exposure
duration can also be obtained as
$D_s=(\Delta t/R_{\max})\Sigma_s$ in person-s. With
$\Delta t=0.1$\,s and $R_{\max}=10$, one reported unit of $\Sigma_s$
corresponds to $0.1$ risk-score--person--s or $0.01$
full-intensity-equivalent person-s. We report the raw accumulator because
it is the direct simulation output and $\Delta t$ and $R_{\max}$ are fixed
in every comparison.

$\Sigma_s$, hereafter Soft Risk, measures how long and how strongly the simulated population
occupies high-score cells. It is a relative model-exposure proxy. Because
$R$ has not been calibrated against a physical concentration, thermal
dose, FED, or dose--response relation, neither $E_i$ nor $\Sigma_s$ is a
physical dose or an estimate of injury probability.

The Hard Harm tally is
\begin{equation}
H = \sum_n \sum_{i\in\mathcal{A}_n}
\mathbb{I}\!\left[
R\!\big(\mathbf{x}_i(t_n),t_n\big)>R_H
\right],
\label{eq:threshold_exceedance}
\end{equation}
where $R_H=4.5$ is the model-defined Hard Harm threshold. Thus,
$H$ counts agent--time-step pairs above $R_H$, rather than distinct
agents. All simulations use $\Delta t=0.1$\,s; the corresponding
aggregate above-threshold occupancy duration is $\Delta t\,H$ person-s. Because
$\Delta t$ is fixed across all configurations, the reported raw tally
$H$ preserves the same comparisons while retaining the integer-valued
simulation output. Accordingly, $H$ can exceed the number of agents in
a run and must not be interpreted as an affected-person count.

Hard Harm $H$ is also distinct from the blocking mask $B(\mathbf{p},t)$.
The mask closes navigation cells at $R_{\mathrm{block}}=6.0$, whereas
$H$ counts active-agent samples above $R_H=4.5$. Thus $B$ is a route
feasibility field and $H$ is an outcome accumulator; the two must not be
interchanged even though both use hard thresholds.

For an ensemble of $K$ runs, the cross-run mean is denoted $\mu_K[H]$
according to Eq.~\eqref{eq:ensemble_summaries}.
Under the hazard-intensity model in
Eq.~\eqref{eq:hazard_intensity} with $R_{\max}=10$,
$\alpha_{\mathrm{pre}}=0.30$, and $\tau_{\mathrm{rise}}=120$\,s, the
threshold $R_H=4.5$ corresponds to approximately $29$\,s after local
hazard arrival under the baseline intensity curve.
Exceedance is assumed to indicate elevated potential for adverse
consequences under the simulated core hazard. Because $R$ is not calibrated
to physical dose--response data, however, Hard Harm is interpreted as a
comparative threshold-event proxy rather than a casualty or injury count,
and it does not estimate incapacitation or mortality probability.
The hazard-field parameterisation and field-level response illustrations
are reported in
\appref{sm:hazard_sensitivity}.

% ---------------------------------------------------------
\subsection{Congestion}
\label{sec:eval_congestion}

Congestion is interpreted through the accumulated spatial fields defined
in \appref{sm:spatial_diagnostics}. The main analysis uses cumulative
density load and slow-and-crowded duration to describe congestion magnitude,
spatial extent, and persistence; quantities used only in the update-period
sweep are defined where they are reported.

% ---------------------------------------------------------
\subsection{Reliability and tail risk}
\label{sec:eval_reliability}

Maze-like settlements are prone to tail risks driven by cascading
congestion and evolving hazards. We adopt two reliability-oriented
formulations.

Given engineering thresholds $\tau$ (time) and $\xi$ (exposure):
\begin{equation}
\mathcal{R}_T(\tau) = \mathbb{P}(T_{0.95} \le \tau),
\qquad
\mathcal{R}_E(\xi) = \mathbb{P}(\Sigma_s \le \xi).
\label{eq:reliability}
\end{equation}
Exceedance probabilities are $1 - \mathcal{R}_T(\tau)$ and
$1 - \mathcal{R}_E(\xi)$.

For a non-negative loss $Z$ (e.g.\ $Z = T_{0.95}$ or $\Sigma_s$),
\begin{equation}
\operatorname{VaR}_\eta(Z) = \inf\{z : \mathbb{P}(Z \le z) \ge \eta\},
\qquad
\operatorname{CVaR}_\eta(Z)
= \min_{\zeta \in \mathbb{R}}
\left\{
\zeta + \frac{1}{1-\eta}\,
\mathbb{E}\!\left[(Z-\zeta)_+\right]
\right\}
\label{eq:cvar}
\end{equation}
where $(x)_+=\max(x,0)$~\citep{Rockafellar2002CVaR}. This definition is
valid for continuous and discrete loss distributions, including cases
with probability mass at the VaR threshold; the often-used conditional
mean $\mathbb{E}[Z\mid Z\geq\operatorname{VaR}_\eta(Z)]$ is not generally
equivalent in such cases. In a probabilistic reliability analysis,
$\operatorname{CVaR}_{0.95}$ therefore summarises the upper loss tail.
This explicit tail summary extends evacuation studies that evaluate
dynamic guidance using cumulative risk, exposure, or stochastic
disruption outcomes rather than mean travel time alone
~\citep{He2024DynamicMultiHazardEvac,Luan2025ToxicGasAirport,Yang2025MetroStampede}.


% ---------------------------------------------------------
\subsection{Assessment protocol}
\label{sec:assessment_protocol}

All navigation strategies are executed with matched random seeds at
each scenario point, enabling strictly paired comparisons. Superscript
$(k)$ denotes a run-level quantity: for example,
$T_p^{(k)}$ is a within-run agent quantile, whereas
$\Sigma_s^{(k)}$ and $H^{(k)}$ are run-level population aggregates. To
distinguish these quantities from summaries across runs, we use
\begin{equation}
\mu_K[X] = \frac{1}{K}\sum_{k=1}^{K}X^{(k)},
\qquad
Q_{q,K}[X] = \inf\!\left\{x:
\frac{1}{K}\sum_{k=1}^{K}\mathbb{I}\!\left[X^{(k)}\le x\right]\ge q
\right\}.
\label{eq:ensemble_summaries}
\end{equation}
Here $\mu_K[\cdot]$ is the cross-run mean and $Q_{q,K}[\cdot]$ the
cross-run empirical $q$-quantile. Thus, the subscript on $T_p^{(k)}$
always refers to an agent-level quantile within run $k$;
the subscript $q,K$ on $Q_{q,K}$ refers to aggregation across runs.
The $\beta$ sweep uses $Q_{0.95,K}$ as a conservative operating-point
criterion under adverse scenario realisations, whereas the update-period
sweep and Shipai evaluation use $\mu_K$ to quantify expected performance
and computation across their ensembles. These functionals answer
different questions and are therefore labelled explicitly rather than
treated as directly comparable estimates.

When a
probabilistic scenario ensemble is available, $\mathcal{R}_T$,
$\mathcal{R}_E$, and CVaR are estimated over $K$ matched trials. The
uncertain inputs may include demand level, initial spatial
distribution, hazard origin, and hazard-growth parameters, depending
on the experiment. Empirical reliability estimates are
\begin{equation}
\widehat{\mathcal{R}}_T(\tau)
= \frac{1}{K} \sum_{k=1}^{K}
  \mathbb{I}\!\big(T_{0.95}^{(k)} \le \tau\big),
\label{eq:reliability_mc}
\end{equation}
and, for trial losses $Z_1,\ldots,Z_K$, the empirical tail estimate is
\begin{equation}
\widehat{\operatorname{CVaR}}_\eta
= \min_{\zeta \in \mathbb{R}}
\left\{
\zeta + \frac{1}{(1-\eta)K}
\sum_{k=1}^{K}(Z_k-\zeta)_+
\right\}.
\label{eq:cvar_mc}
\end{equation}
When $(1-\eta)K$ is an integer, this equals the mean of the largest
$(1-\eta)K$ losses; otherwise, the minimisation gives the appropriate
fractional weight at the empirical quantile. We use bootstrap resampling
to construct $95\%$ confidence intervals ($B = 1{,}000$ resamples
unless otherwise stated). For a paired contrast between strategies
$a$ and $b$, the scenario-level difference is
$d^{(k)}=X_a^{(k)}-X_b^{(k)}$; bootstrap samples are drawn from the
$K$ matched scenario indices, and the mean difference and percentile
interval are computed from the resampled pairs. An interval containing
zero is treated as an absence of evidence for a mean difference, not as
evidence that the strategies are equivalent. The Shipai summaries and
tail-time contrasts use $B=2{,}000$ resamples. The main reported results
emphasise paired strategy comparisons and high-quantile outcomes;
\appref{sm:disaster_locations} illustrates how spatial
uncertainty in hazard origin affects evacuation tail dispersion.


% =======================================================
% Section: Model Calibration and Validation
% =======================================================
\FloatBarrier
\section{Model Calibration and Validation}
\label{sec:calibration_validation}

Controlled evacuation experiments at the community scale are constrained
by ethics, safety, and instrumentation limits. We therefore adopt a
\textbf{bottom-up hierarchical validation strategy}: if the model
reproduces fundamental pedestrian relations at the micro scale and
emergent bottleneck behaviour at the meso scale, its macro-scale use
for system-level analysis becomes
defensible~\citep{seyfried2005fundamental,zhang2013empirical}. This
section reports (i) calibration and quantitative error analysis against
empirical fundamental diagrams (\S\ref{sec:calibration_fd}), and
(ii) validation of crowd dynamics at representative geometric
discontinuities (\S\ref{sec:validation_flow_width}).


\subsection{Micro-level calibration: fundamental diagram}
\label{sec:calibration_fd}

The model is calibrated against the pedestrian \emph{fundamental
diagram}---the joint relationship among density ($\rho$), speed ($v$),
and flow ($q = \rho v$)---which constitutes a canonical benchmark for
unidirectional pedestrian
flow~\citep{seyfried2005fundamental,zhang2013empirical}. Two features
are robustly observed across experimental studies: speed decreases
monotonically with density, and flow exhibits a unimodal profile that
attains its maximum at a critical density. Reproducing both is a
necessary condition for any agent-based model intended for quantitative
evacuation analysis.

\begin{figure}[htbp]
    \centering
    \includegraphics[width=1.0\linewidth]{figures/05_model_calibration_validation/paradigm_obs_mixed_styled.png}
    \caption{Empirical speed--density (left) and flow--density (right)
    relationships from corridor experiments~\citep{zhang2013empirical}.
    The sample cloud is obtained by mixed sampling from eight independent
    evacuation trajectory files, corresponding to eight crowd experiments;
    colours indicate the source file and the black curves show binned
    medians. The aggregated relationships are used here as the calibration
    target.}
    \label{fig:fundamental_diagram_obs}
\end{figure}

The procedure has two stages.

\emph{Stage 1: closure fit.} The speed--density closure introduced in
\S\ref{sec:speed_density_cost} (Eq.~\ref{eq:speed_field}) is fitted to
the paired observations $\{(\rho_j, v_j)\}$ reported by
\citet{zhang2013empirical} (Fig.~\ref{fig:fundamental_diagram_obs}) by
constrained nonlinear least squares (with $v_f, \rho_0, n > 0$),
yielding $\hat{v}_f = 1.574$\,m/s,
$\hat{\rho}_0 = 1.735$\,ped/m$^{2}$, and $\hat{n} = 0.852$ with a
residual RMSE of $0.037$\,m/s on the empirical sample.

\emph{Stage 2: model-level reproduction.} The calibrated closure is
embedded in the full agent-based model and evaluated through controlled
corridor simulations designed to emulate the narrow ``handshake''
alleys characteristic of urban informal settlements. Agent counts are
incremented systematically so that the realised density spans the
free-flow, transitional, and congested regimes, and instantaneous
local density and velocity are extracted from the simulated
trajectories for comparison with the empirical reference. The
speed--density relation is reproduced with $R^{2} = 0.943$ and the
overall flow--density trend with $R^{2} = 0.810$; the lower $R^{2}$ on
flow is expected because $q = \rho v$ compounds independent estimation
errors in $\rho$ and $v$. Discrepancies are concentrated in the
high-density regime ($\rho > 4$\,ped/m$^{2}$), where flow is governed
by local fluctuations and stop-and-go instabilities; we treat this
regime as a known limitation of the present calibration and return to
it when interpreting the system-level results in
\S\ref{sec:shipai}.


% =========================================================
\subsection{Meso-level validation: geometric flow capacity}
\label{sec:validation_flow_width}

Micro-level calibration (\S\ref{sec:calibration_fd}) confirms that
individual agents reproduce the local speed--density relationship.
However, reproducing local kinematics does not guarantee that collective
flow responds correctly to geometric discontinuities---the dominant
bottleneck type in dense urban informal settlements. We therefore test
whether the calibrated model also captures capacity changes induced by
corridor width, turns, and merges.

Three controlled configurations isolate these effects:
straight corridors, L-shaped corners ($90^\circ$ turns), and T-junctions
(Fig.~\ref{fig:corridor_geometries}). For each configuration, effective
walking width is varied from 1.0 to 5.0\,m in 1.0\,m increments, with
five independent runs per width, giving 75 scheduled runs. The response
variable is the maximum sustained flow rate $J_{\max}$, computed over
the steady evacuation phase ($5\%$--$90\%$ completion) to exclude
start-up transients and tail effects. Before fitting, runs were screened
within each geometry--width group using a common completion-time rule:
a run was flagged when its final evacuation time exceeded
$Q_3+1.5\,\mathrm{IQR}$ (or was missing). This flagged three high-duration
runs---the L-corner at $b=3$\,m ($t_{\mathrm{end}}=235$\,s) and the
T-junction at $b=2$ and $5$\,m ($252$ and $4{,}999$\,s, respectively)---so
the screened fit uses 72 runs. Because this filter is outcome-based, we
also report the unfiltered sensitivity below rather than treating the
flagged cases as measurement errors.

\begin{figure*}[htbp]
    \centering
    \includegraphics[width=\textwidth]{figures/05_model_calibration_validation/corridor_geometries.pdf}
    \caption{Canonical corridor configurations used for meso-level
    validation: (a) straight, (b) L-shaped, and (c) T-shaped, isolating
    the effects of width, turning, and merging on flow capacity.}
    \label{fig:corridor_geometries}
\end{figure*}

For the straight corridor, the empirical benchmark is the flow--width relation
\begin{equation}
J_{\max} = C \cdot b,
\label{eq:flow_width}
\end{equation}
where $b$ is corridor width and $C$ is specific capacity. Controlled
straight-corridor experiments report $J\simeq1.6b$, while the broader
published range is approximately $C=1.6$--$2.2$\,person/(m$\cdot$s)
\citep{weidmann1993fundamental,seyfried2005fundamental,zhang2013empirical}.
This provides a direct quantitative comparator for the straight case.
The available corner and T-junction evidence is informative but not
geometry-matched capacity--width validation. In a fixed 2.4-m setup,
the fundamental diagrams measured immediately before and after a
$90^\circ$ corner largely overlap, whereas the T-junction branches show
lower pre-merge velocities and a specific-flow plateau relative to the
post-merge stream~\citep{zhang2012turning,zhang2013tjunction}. A separate
bend experiment also finds spatially nonuniform velocity, with the
largest reduction near the inner corner~\citep{dias2022bend}. These
studies therefore supply qualitative geometric expectations, but not
capacity--width bands for the L-corner or T-junction.

Across the three simulated geometries, fitted flow scales approximately
linearly with width ($R^{2}\ge0.978$; Fig.~\ref{fig:flow_width_validation}).
The straight-corridor estimate ($C_{\mathrm{S}}=1.90$) lies within the
external range. The internal ratios
$C_{\mathrm{L}}/C_{\mathrm{S}}\approx1.05$ and
$C_{\mathrm{T,arm}}/C_{\mathrm{S}}\approx0.61$ indicate similar aggregate
capacity for the straight and L-shaped cases and lower per-arm throughput
under symmetric merging. The conclusion is insensitive to the completion-time
screen: using all 75 scheduled runs changes the fitted total slopes
$(C_{\mathrm{S}},C_{\mathrm{L}},C_{\mathrm{T,total}})$ from
$(1.90,1.98,2.30)$ to $(1.90,1.96,2.29)$ and changes the corresponding
$R^{2}$ values from $(0.978,0.993,0.997)$ to
$(0.978,0.987,0.998)$. These are internal geometric sensitivity results.
Direct validation of bend and T-junction capacity would require matched
geometries, demand conditions, measurement regions, and trajectory-level
spatial statistics.

\begin{table}[htbp]
    \centering
    \caption{Fitted simulation capacities from the 72-run screened
    analysis and the scope of external comparison. Only the straight
    corridor has a directly comparable empirical capacity range; the
    bend and T-junction sources provide qualitative spatial-flow evidence
    under fixed, different geometries.}
    \label{tab:flow_width_val}
    \small
    \begin{tabularx}{\textwidth}{@{}
      >{\raggedright\arraybackslash}p{0.18\textwidth}
      >{\centering\arraybackslash}p{0.08\textwidth}
      >{\centering\arraybackslash}p{0.08\textwidth}
      >{\raggedright\arraybackslash}p{0.28\textwidth}
      >{\raggedright\arraybackslash}X@{}}
    \toprule
    \textbf{Scenario} & $C^{\mathrm{sim}}$ & $R^{2}$ &
        \textbf{External comparator} & \textbf{Interpretation} \\
    \midrule
    Straight corridor & 1.90 & 0.978 & 1.6--2.2 person/(m$\cdot$s)
      & Quantitative agreement with the cited range. \\
    L-corner ($90^\circ$) & 1.98 & 0.993 & Before/after-corner
      fundamental diagrams and spatial velocity variation;
      no comparable capacity--width range
      & Internal width-scaling check only; local bend effects are not
        validated by this aggregate statistic. \\
    T-junction (per arm)$^{\dagger}$ & 1.15 & 0.997 & Pre-merge velocity reduction and
      branch-flow plateau; no comparable per-arm capacity range
      & Qualitative consistency with merge restriction; magnitude is not
        externally validated. \\
    \bottomrule
    \end{tabularx}
    \vspace{2pt}
    \parbox{\textwidth}{\footnotesize ${}^{\dagger}$The plotted T-junction
    fit uses total two-branch inflow; the table reports half of that slope
    as a symmetric per-arm equivalent.}
\end{table}

\begin{figure}[htbp]
    \centering
    \includegraphics[width=0.92\textwidth]{figures/05_model_calibration_validation/flow_width_validation.pdf}
    \caption{Three-way comparison of simulated flow--width relationships.
    Faint symbols denote retained runs, large symbols and error bars show
    the per-width mean $\pm$ one standard deviation, and solid coloured
    lines are origin-constrained OLS fits. Grey crosses identify the three
    high-duration runs flagged by the stated completion-time rule. The blue
    band is the published straight-corridor capacity range
    $C=1.6$--$2.2$\,person/(m$\cdot$s); no empirical band is assigned to
    the fixed-width L-corner or T-junction studies. The T-junction series
    is total two-branch inflow, whereas Table~\ref{tab:flow_width_val}
    reports its symmetric per-arm equivalent.}
    \label{fig:flow_width_validation}
\end{figure}







% =========================================================
\section{Strategy Calibration on a Simplified Network}
\label{sec:strategy_calibration}

The proposed Density--Hazard strategy involves two parameters that
govern its operating behaviour: the hazard-sensitivity weight $\beta$
in Eq.~\eqref{eq:phi_continuous}, which sets the relative strength of
the hazard term against the density term, and the field update period
$\Delta t_u$, which determines the temporal resolution at which the
guidance field tracks the evolving crowd state. Both must be specified
before the strategy is applied to the full-scale case in
\S\ref{sec:shipai}.

\begin{figure}[htbp]
    \centering
    \includegraphics[width=0.55\linewidth]{figures/06_strategy_calibration_simplified_network/simplenetwork_2.pdf}
    \caption{Simplified calibration network used for the strategy
    sensitivity experiments. Grey polygons denote built-up or
    non-walkable blocks, white corridors denote the pedestrian movement
    network, and circled green markers denote boundary exits. The
    reduced geometry preserves narrow alleys, irregular junctions,
    competing route families, and multiple evacuation destinations while
    substantially reducing simulation cost relative to the full Shipai
    benchmark.}%; the scale bar indicates 50 m.}
    \label{fig:simplified_network}
\end{figure}

To isolate parameter effects from scene-specific confounds and to enable
the Monte Carlo replicates required for stable bootstrap estimates, all
calibration experiments in this section are conducted on a
\emph{simplified network} that retains the principal topological
features of urban-village evacuation---narrow corridors, multiple
branching points, and competing path families---while reducing
computational cost by an order of magnitude relative to the Shipai
benchmark. The update-period calibration uses $K = 20$ independent
scenario instances per configuration, while the strategy and
$\beta$-sensitivity sweep uses a paired seven-strategy set of $50$
instances and reports the $K=49$ runs retained after applying the
prespecified IQR filter. All confidence intervals are $95\%$ bootstrap
with $B = 1{,}000$. The transferability of the resulting operating points
$(\beta^{*}, \Delta t_u^{*})$ to the full Shipai network is examined in
\S\ref{sec:shipai}.




The three strategies used in the main comparisons are all special cases
of the unified framework introduced in \S\ref{sec:method}, distinguished
by their parameter configuration at the selected operating point
(Table~\ref{tab:strategy_mapping}). The calibration sweeps that vary
$\beta$ or $\Delta t_u$ around this point are reported separately below.

\begin{table}[htbp]
\centering
\small
\setlength{\tabcolsep}{5pt}
\caption{Main comparison strategies at the selected operating point,
expressed as complete parameter configurations of the unified framework
(Eq.~\ref{eq:phi_continuous}).}
\label{tab:strategy_mapping}
\begin{tabular}{lccccp{0.32\linewidth}}
\toprule
\textbf{Strategy} & $\alpha$ & $\beta$ (s/m) & $\gamma$ & $\Delta t_u$ & \textbf{Description} \\
\midrule
Static     & 1 & 0 & 0   & $\infty$    & Pure geometric distance-to-exit field, computed once; no density or hazard input. \\
Density-60 & 1 & 0 & 0   & \SI{60}{s} & Density-dependent travel-time cost only; neither $R$ nor $B$ enters navigation. \\
DH-60      & 1 & 1 & $+\infty$ & \SI{60}{s} & Density-dependent travel-time cost, continuous hazard cost, and hard blocking. \\
\bottomrule
\end{tabular}
\end{table}

For Static, the edge cost is defined as
$c_{\mathrm{Static}}(\mathbf{p}\to\mathbf{q})=
c_0\Delta s_{\mathbf{pq}}$, with the spatially uniform reference cost
$c_0=\SI{1}{\second\per\meter}$. Because $c_0$ is constant, this is
exactly a geometric shortest-distance field up to a common scale. Its
potential is computed once from the walkable geometry.
It does not evaluate $c_p$ from the initial population and is
therefore neither an initial-density-frozen field nor an
initial-Density--Hazard-frozen field. The value $\alpha=1$ in
Table~\ref{tab:strategy_mapping} is only the common scale convention;
it does not introduce density dependence into Static.

Density-60 instead recomputes the travel-time field $c_p$ from the
instantaneous crowd density every \SI{60}{s}, while retaining
$\beta=\gamma=0$. DH-60 uses the same density-dependent travel-time
channel and update period, and additionally uses the continuous hazard
term ($\beta=\SI{1}{\second\per\meter}$) and exact blocking
($\gamma=+\infty$). These definitions give the intended nested
comparison: geometric distance only, density-dependent travel time, and
density-dependent travel time plus hazard.

The hazard process is simulated under all three strategies for paired
outcome evaluation. However, it is observational only for Static and
Density-60: agents accumulate exposure according to the hazard field
along their realised trajectories, but that field cannot alter their
route choice. Thus Density-60 has no access to either continuous hazard
intensity $R$ or the derived blocking mask $B$.


\subsection{Hazard-weight calibration}
\label{sec:beta_calibration}

We sweep $\beta \in \{0.1, 0.5, 1.0, 5.0, 10.0\}$\,s/m with
$\alpha = 1.0$ and exact hard blocking ($\gamma=+\infty$) fixed, and
report four
lower-is-better indicators:
the cross-run upper quantiles $Q_{0.95,K}[T_{0.95}]$ and
$Q_{0.95,K}[T_{0.99}]$
(\S\ref{sec:eval_efficiency}), the cross-run upper population-exposure summary
$Q_{0.95,K}[\Sigma_s]$
(\S\ref{sec:eval_exposure}, Eq.~\ref{eq:soft_risk_population}), and the cross-run
mean Hard Harm tally $\mu_K[H]$ at threshold $R_H=4.5$
(\S\ref{sec:eval_exposure}, Eq.~\ref{eq:threshold_exceedance}).


\begin{table}[htbp]
    \centering
    \caption{$\beta$-sensitivity of Density--Hazard. Values for
             $T_{0.95}^{(k)}$, $T_{0.99}^{(k)}$, and $\Sigma_s^{(k)}$
             are summarised by the cross-run 95th quantile
             $Q_{0.95,K}[\cdot]$ with 95\% bootstrap CIs; $H^{(k)}$ is
             summarised by the cross-run mean $\mu_K[H]$. Hard Harm is a
             threshold-event tally in agent--time-step pairs, whereas Soft Risk
             $\Sigma_s$ is reported in $10^6$
             risk-score--person--time-step units. The paired sweep uses $K=49$ runs after
             the prespecified IQR filter; the highlighted $\beta = 1.0$ row
             is the recommended operating point.}
    \label{tab:beta_sensitivity}
    \footnotesize
    \setlength{\tabcolsep}{3.5pt}
    \renewcommand{\arraystretch}{1.15}
    \begin{tabular}{ccccc}
    \toprule
    $\boldsymbol{\beta}$ (s/m)
        & \makecell{$Q_{0.95,K}[T_{0.95}]$\\(s)}
        & \makecell{$Q_{0.95,K}[T_{0.99}]$\\(s)}
        & \makecell{$Q_{0.95,K}[\Sigma_s]$\\($\times10^{6}$ RSU)\\(person--step)}
        & $\mu_K[H]$ \\
    \midrule
    0.1
        & \makecell{256.4 \\ {\scriptsize [234.0, 317.4]}}
        & \makecell{311.7 \\ {\scriptsize [280.1, 364.8]}}
        & \makecell{1.95 \\ {\scriptsize [1.72, 3.24]}}
        & \makecell{1758.7 \\ {\scriptsize [244.3, 4285.1]}} \\
    0.5
        & \makecell{265.3 \\ {\scriptsize [249.6, 323.8]}}
        & \makecell{321.1 \\ {\scriptsize [301.6, 379.8]}}
        & \makecell{1.67 \\ {\scriptsize [1.44, 2.31]}}
        & \makecell{110.4 \\ {\scriptsize [0.0, 331.1]}} \\
    \rowcolor{gray!12}
    \textbf{1.0}
        & \makecell{277.2 \\ {\scriptsize [253.9, 331.9]}}
        & \makecell{327.6 \\ {\scriptsize [310.1, 384.7]}}
        & \makecell{1.55 \\ {\scriptsize [1.37, 1.98]}}
        & \makecell{10.9 \\ {\scriptsize [0.0, 32.8]}} \\
    5.0
        & \makecell{345.3 \\ {\scriptsize [278.1, 382.8]}}
        & \makecell{420.3 \\ {\scriptsize [338.8, 462.0]}}
        & \makecell{1.42 \\ {\scriptsize [1.20, 1.55]}}
        & \makecell{0.0 \\ {\scriptsize [0.0, 0.0]}} \\
    10.0
        & \makecell{353.9 \\ {\scriptsize [283.4, 404.2]}}
        & \makecell{434.0 \\ {\scriptsize [342.7, 488.9]}}
        & \makecell{1.38 \\ {\scriptsize [1.20, 1.46]}}
        & \makecell{0.0 \\ {\scriptsize [0.0, 0.0]}} \\
    \bottomrule
    \end{tabular}
\end{table}


The sensitivity sweep is summarised in
Table~\ref{tab:beta_sensitivity} and
Fig.~\ref{fig:beta_tradeoff_combined}. This sweep uses the paired
seven-strategy ensemble after applying the prespecified IQR filter
($K=49$). As $\beta$ increases from $0.1$
to $10.0$, $Q_{0.95,K}[T_{0.95}]$ rises monotonically by $38.0\%$
($256.4 \to 353.9$\,s) while $Q_{0.95,K}[\Sigma_s]$ falls monotonically by
$29.1\%$ ($1.95 \to 1.38 \times 10^{6}$), confirming that the hazard
term acts in the intended direction throughout the tested range. The
trade-off is strongly nonlinear: increasing $\beta$ from $0.1$ to $1.0$
reduces $Q_{0.95,K}[\Sigma_s]$ by $20.3\%$ at an $8.1\%$ time cost, whereas
increasing from $1.0$ to $10.0$ yields a further $11.1\%$ exposure
reduction but raises $Q_{0.95,K}[T_{0.95}]$ by $27.7\%$.


\begin{figure}[htbp]
    \centering
    \begin{minipage}[t]{0.49\linewidth}
        \centering
        \includegraphics[width=\linewidth]{figures/06_strategy_calibration_simplified_network/4x2_beta_tradeoff.pdf}
        \small (a) $Q_{0.95,K}[T_{0.95}]$--$Q_{0.95,K}[\Sigma_s]$ trade-off.
    \end{minipage}
    \hfill
    \begin{minipage}[t]{0.49\linewidth}
        \centering
        \includegraphics[width=\linewidth]{figures/06_strategy_calibration_simplified_network/4x2_beta_tradeoff_harm.pdf}
        \small (b) $Q_{0.95,K}[T_{0.95}]$--$\mu_K[H]$ trade-off.
    \end{minipage}
    \caption{Trade-off analysis across $\beta$ (in s/m). Left: monotonic
    efficiency--exposure trade-off, with $\beta = 1.0$ marking a knee
    before the high-$\beta$ time penalty dominates. Right: Hard Harm
    decreases with $\beta$ and is eliminated at $\beta \ge 5.0$, while
    $\beta = 1.0$ nearly eliminates Hard Harm without the high-$\beta$
    detour cost. Static (square) and Density-aware (diamond) are shown
    for reference; points show filtered paired estimates ($K=49$).}
    \label{fig:beta_tradeoff_combined}
\end{figure}

\begin{table}[htbp]
    \centering
    \caption{Trend summary across $\beta$ in the paired sensitivity
    sweep. Lower is better on every metric.}
    \label{tab:ranking_stability}
    \small
    \renewcommand{\arraystretch}{1.18}
    \begin{tabular}{lll}
    \toprule
    \textbf{Metric} & \textbf{Observed pattern} & \textbf{Implication} \\
    \midrule
    $Q_{0.95,K}[T_{0.95}]$ & increases with $\beta$ & high $\beta$ is time-costly \\
    $Q_{0.95,K}[\Sigma_s]$ & decreases with $\beta$ & larger $\beta$ lowers exposure \\
    $\mu_K[H]$  & falls to zero by $\beta \ge 5.0$ & $\beta=1.0$ is a knee point \\
    \bottomrule
    \end{tabular}
\end{table}


The Hard Harm metric decreases monotonically over the tested range.
$\mu_K[H]$ falls from $1758.7$ at $\beta = 0.1$ to $110.4$ at $\beta = 0.5$ and $10.9$ at
$\beta = 1.0$, and is eliminated under the same IQR-filtered paired
sample at $\beta \in \{5.0, 10.0\}$. Thus $\beta = 1.0$ is not a
Hard Harm global minimum; it is the knee of the efficiency--exposure
curve. Larger $\beta$ values remove the remaining Hard Harm events but
do so by inducing much longer detours, increasing $Q_{0.95,K}[T_{0.95}]$ by
$24.6\%$--$27.7\%$ relative to $\beta = 1.0$.

The ranking pattern is therefore metric-dependent
(Table~\ref{tab:ranking_stability}). Combining the exposure reduction,
near-elimination of Hard Harm, and avoidance of the high-$\beta$
tail-time penalty, we adopt $\beta^{*} = 1.0$ as the recommended
operating point.

\subsection{Field-update-period calibration}
\label{sec:update_period}

We calibrate the field update period $\Delta t_u$ using the
Density-aware strategy with $\beta=\gamma=0$, thereby isolating the
temporal responsiveness of the density-dependent travel-time field from
hazard avoidance. Smaller $\Delta t_u$ yields a more responsive field at
higher computational cost; larger $\Delta t_u$ amortises recomputation
but allows the field to lag behind fast-evolving crowd states. We compare
the separately constructed, pure-distance Static reference against five
Density-aware settings with
$\Delta t_u \in \{10, 30, 60, 120, 180\}$\,s. Static is not a
Density--Hazard field frozen at $t=0$, nor is its cost derived from the
initial crowd density. After selecting $\Delta t_u^{*}$ in this isolated
density-feedback experiment, we use the same period for DH-60 so that
the two dynamic strategies differ only in their use of $R$ and $B$.

This sweep uses $K=20$, and every plotted point is a cross-run mean
$\mu_{20}[\cdot]$. Here $T_p^{(k)}$ is the within-run clearance
quantile, $M_{8}^{+,(k)}$ and $N_{8}^{+,(k)}$ are the run-level
Layer-8 summaries, and $C_{\mathrm{tot}}^{(k)}$ and
$C_{\mathrm{field}}^{(k)}$ are the run-level computation times.

The congestion diagnostics in this sweep are derived from the cumulative
Layer-8 field. We report its positive-cell mean, $M^{+}_{8}$, and the number
of cells with nonzero accumulation, $N^{+}_{8}$, as measures of congestion
magnitude and spatial extent, respectively.


\begin{figure}[htbp]
    \centering
    \includegraphics[width=\linewidth]{figures/06_strategy_calibration_simplified_network/fig3_tradeoff.pdf}
    \caption{Pareto trade-off between evacuation quality and
    computation cost for the pure-distance Static reference and five
    Density-aware update periods. (a) $\mu_{20}[T_{0.95}]$ vs.\
    $\mu_{20}[C_{\mathrm{tot}}]$; (b) $\mu_{20}[M^{+}_{8}]$ vs.\
    $\mu_{20}[C_{\mathrm{field}}]$.
    Dashed curves connect dynamic
    configurations in order of increasing $\Delta t_u$; static
    baseline shown as a filled square. Error bars: 95\% bootstrap CI.}
    \label{fig:tradeoff}
\end{figure}

All dynamic configurations improve evacuation performance over the
static baseline ($\mu_{20}[T_{0.95}] = \SI{422.5}{s}$,
$\mu_{20}[C_{\mathrm{tot}}]=\SI{18.0}{s}$ per run). The most aggressive setting
($\Delta t_u = \SI{10}{s}$) attains the largest gains---$-16.5\%$ in
$\mu_{20}[T_{0.95}]$ ($\to \SI{352.6}{s}$), $-20.0\%$ in
$\mu_{20}[T_{0.99}]$, and $-24.6\%$ in $\mu_{20}[M^{+}_{8}]$---but incurs a $5.59\times$
increase in total computation ($\SI{100.5}{s}$, of which $85.2\%$ is
field recomputation; Fig.~\ref{fig:tradeoff}). As $\Delta t_u$
relaxes, the tail-time gain relative to Static shrinks while computational
cost decreases along the trade-off curve.


\begin{figure}[htbp]
    \centering
    \includegraphics[width=\linewidth]{figures/06_strategy_calibration_simplified_network/fig4_marginal_benefit.pdf}
    \caption{Sensitivity of four performance metrics to $\Delta t_u$
    (log-scaled axis): (a) $\mu_{20}[T_{0.95}]$, (b)
    $\mu_{20}[T_{0.99}]$, (c) $\mu_{20}[M^{+}_{8}]$, and (d)
    $\mu_{20}[N^{+}_{8}]$. Shaded bands: 95\% bootstrap CI; horizontal
    dashed: static baseline. The metric-specific change-point markers are
    located at \SI{60}{s} in panels (a)--(b) and \SI{30}{s} in panels
    (c)--(d).}
    \label{fig:diminishing}
\end{figure}



Figure~\ref{fig:diminishing} shows that the apparent elbow is
metric-dependent: the change point occurs at \SI{60}{s} for
$\mu_{20}[T_{0.95}]$ and $\mu_{20}[T_{0.99}]$, but at \SI{30}{s} for
$\mu_{20}[M^{+}_{8}]$ and $\mu_{20}[N^{+}_{8}]$.
Accordingly, \SI{60}{s} is not interpreted as a universal optimum across
all diagnostics. It is selected from the joint tail-time--computation
trade-off: relative to \SI{10}{s}, it increases $\mu_{20}[T_{0.95}]$ by only
\SI{15.0}{s} while reducing total runtime by a factor of $3.4$; relaxing
further from \SI{60}{s} to \SI{120}{s} saves only \SI{6.8}{s} of runtime
while increasing $\mu_{20}[T_{0.95}]$ by \SI{12.0}{s}. The selected period
therefore represents a balanced operating point rather than a
single-metric optimum.



\begin{figure}[htbp]
    \centering
    \includegraphics[width=0.75\linewidth]{figures/06_strategy_calibration_simplified_network/fig5_compute_vs_period.pdf}
    \caption{Cross-run-mean computation-time decomposition vs.\ $\Delta t_u$.
    Circles: $\mu_{20}[C_{\mathrm{tot}}]$; squares:
    $\mu_{20}[C_{\mathrm{field}}]$;
    triangles (right axis): dynamic-field fraction of total runtime.
    Error bands: 95\% bootstrap CI; horizontal axis log-scaled.}
    \label{fig:compute_cost}
\end{figure}

Fig.~\ref{fig:compute_cost} decomposes total runtime into the
field-update component and the remainder. Agent-state updates are
GPU-accelerated on an NVIDIA GeForce RTX 5090, whereas the online
navigation-field recomputation is still performed on the CPU (Intel Core
Ultra 9 285K, 24 logical processors). The field-update component in
Fig.~\ref{fig:compute_cost} therefore quantifies this CPU-side
bottleneck: its runtime fraction falls from $85.2\%$ at
$\Delta t_u = \SI{10}{s}$ to $49.9\%$ at $\SI{60}{s}$ ($\SI{15.0}{s}$
of $\SI{30.0}{s}$ total) and converges to the static baseline beyond
\SI{120}{s}. The chosen operating point
$\Delta t_u^{*} = \SI{60}{s}$ thus delivers a $5.7\times$ reduction in
field-update cost relative to $\SI{10}{s}$ for only a \SI{15.0}{s}
increase in $\mu_{20}[T_{0.95}]$, and is adopted as the default for all
subsequent experiments.

Taken together, the simplified-network experiments serve as a
calibration screen for the two operating choices that cannot be fixed
from first principles. The $\beta$ sweep establishes the
efficiency--exposure shape of hazard aversion: increasing $\beta$
consistently lowers Soft Risk and Hard Harm, but values above
$\beta=1.0$ purchase the remaining exposure reductions mainly through a steep
tail-time penalty. Separately, with $\beta=\gamma=0$, the
Density-aware update-period sweep identifies the temporal-resolution
trade-off of density feedback: \SI{10}{s} updates are most responsive but
CPU-bottlenecked, whereas \SI{120}{s}--\SI{180}{s} updates begin to lag
behind transient congestion. The selected operating point
$(\beta^{*}, \Delta t_u^{*})=(1.0,\SI{60}{s})$ therefore represents a
combination of the independently selected hazard weight and update
period rather than a single-metric optimum. It preserves most of the
clearance and risk-control benefits of dynamic guidance while keeping
the online field-update burden within a practical range. The
next section tests whether this simplified-network calibration transfers
to the larger and more heterogeneous Shipai benchmark.


\FloatBarrier
\section{Case Study: Shipai Village}
\label{sec:shipai}
\FloatBarrier


\subsection{Benchmark setting and strategies}
\label{sec:shipai_benchmark}

This section tests whether the operating point
$(\beta^{*}, \Delta t_u^{*})=(1.0,\SI{60}{s})$ identified in
\S\ref{sec:strategy_calibration} transfers from the simplified
calibration network to the full Shipai urban-village benchmark---a
realistic, larger-scale, and topologically heterogeneous setting.

\begin{figure*}[htbp]
  \centering
  \includegraphics[width=0.92\textwidth]{figures/07_shipai_case/shipai_layout.png}
  \caption{Shipai Village benchmark. The layout figure shows the
  remote-sensing base map (a), building footprints (b), and pedestrian
  network (c) used to construct the full-scale evacuation domain. The
  morphology combines dense building blocks, narrow alley corridors, and
  multiple boundary exits, providing a realistic stress test for dynamic
  evacuation guidance in maze-like urban villages.}
  \label{fig:shipai_layout}
\end{figure*}

Shipai was selected as the full-scale benchmark because it combines practical importance and morphological complexity among Guangzhou's high-density urban villages. Located in central Tianhe and surrounded by commercial, educational, medical, and transit facilities, it functions as a mixed residential--employment settlement rather than a closed residential enclave. Local reporting describes a $0.31\,\mathrm{km}^2$ core with more than $20{,}000$ incoming workers and nearly $300$ village lanes, only $25$ of which are wider than $2\,\mathrm{m}$; broader public descriptions of the $0.73\,\mathrm{km}^2$ village report more than $50{,}000$ migrant residents \citep{guangzhouTianhe2023ShipaiMap,pixelNotes2019Shipai}. Because these figures use different spatial boundaries and population definitions, the randomised demand of $40\,087$--$49\,892$ agents (mean $45\,102$) is treated as an upper-load computational stress test rather than a demographic estimate of simultaneous evacuees.



Morphologically, Shipai exhibits a maze-like outdoor pedestrian network of narrow alleys, frequent dead ends, irregular intersections, and strong visibility occlusion (Fig.~\ref{fig:shipai_layout}). Unlike hierarchical street networks, it contains many short links and locally attractive corridor families that become globally inefficient once crowding or hazard propagation alters the available route set: distance-only guidance may overload narrow chokepoints, whereas adaptive rerouting may divert flow into secondary alleys that are safer or more hazardous depending on the evolving field. The benchmark covers a $674\,\mathrm{m} \times 847\,\mathrm{m}$ domain, roughly an order of magnitude larger than the calibration network. We compare the three strategies of Table~\ref{tab:strategy_mapping}---Static (shortest-path), Density-aware with $\Delta t_u^{*} = \SI{60}{s}$ (\emph{Density-60}), and Density--Hazard at the calibrated operating point (\emph{DH-60})---over $K=50$ independent runs per configuration.

\subsection{Aggregate performance metrics}
\label{sec:shipai_metrics}

At deployment scale we retain the same run-level metric definitions but
summarise the $K=50$ runs by $\mu_{50}[\cdot]$, rather than by the
cross-run upper quantile used for the first three metrics in the
$\beta$ calibration. Define, for run $k$,
\begin{equation}
\Sigma_s^{(k)} = \sum_{i \in \mathcal{A}^{(k)}} E_i^{(k)},
\qquad
H^{(k)} = \sum_n\sum_{i\in\mathcal{A}_n^{(k)}}
\mathbb{I}\!\left[
R\!\big(\mathbf{x}_i^{(k)}(t_n),t_n\big)>R_H
\right],
\label{eq:cum_metrics}
\end{equation}
where $\mathcal{A}^{(k)}$ is the agent set in run $k$ and $E_i^{(k)}$ is the
agent-level risk-score accumulator from Eq.~\eqref{eq:exposure};
$\mathcal{A}_n^{(k)}$ is the active-agent set at time step $t_n$.
$\Sigma_s^{(k)}$ is the run-level population Soft Risk, whereas
$H^{(k)}$ is the Hard Harm tally---the cumulative number of
agent--time-step pairs above $R_H=4.5$---consistent
with Eq.~\eqref{eq:threshold_exceedance}. The reported quantities are
$\mu_{50}[\Sigma_s]$ and $\mu_{50}[H]$, with 95\% bootstrap CIs. These
population aggregates are the natural unit for system-level safety
reporting at city-block scale. The use of $\mu_{50}$ targets mean
deployment performance across the matched Shipai scenarios and is kept
distinct from the conservative $Q_{0.95,49}$ criterion used in the
$\beta$ sweep. The separation between Density-60 (for which
$\beta=\gamma=0$ and the hazard field is used only for outcome
evaluation) and DH-60 in
Table~\ref{tab:results_compact} provides the empirical comparison used to
isolate the effect of hazard-aware navigation.

\begin{table*}[htbp]
\footnotesize
\centering
\setlength{\tabcolsep}{4pt}
\renewcommand{\arraystretch}{1.15}
\caption{Performance summary for the Shipai full-scenario benchmark
($K = 50$ runs per strategy). The dynamic settings use the operating
point calibrated on the simplified network and are then evaluated on
the full Shipai domain. Every entry is a cross-run mean
$\mu_{50}[\cdot]$ with a 95\% marginal bootstrap CI;
$\Sigma_s$ (Soft Risk) is reported in $10^6$ risk-score--person--time-step units and
$H$ (Hard Harm) in $10^3$ agent--time-step pairs; $\Delta$ denotes relative improvement over
Static; bold marks the column-best mean. Paired contrasts between
Density-60 and DH-60 are reported in the text.}
\label{tab:results_compact}
\begin{tabular}{lcccc}
\toprule
\textbf{Strategy}
& $\mu_{50}[T_{0.95}]$ (s)
& $\mu_{50}[T_{0.99}]$ (s)
& \makecell{$\mu_{50}[\Sigma_s]$\\($\times10^{6}$ RSU)\\(person--step)}
& $\mu_{50}[H]$ ($\times10^{3}$) \\
\midrule
Static
& \makecell{1114 \\ {\scriptsize [1064, 1169]}}
& \makecell{1431 \\ {\scriptsize [1356, 1517]}}
& \makecell{43.11 \\ {\scriptsize [33.0, 54.7]}}
& \makecell{2649.5 \\ {\scriptsize [1644.8, 3822.7]}} \\
Density-60
& \makecell{\textbf{956} \\ {\scriptsize [931, 983]} \\ {\scriptsize $\Delta = -14.1\%$}}
& \makecell{1228 \\ {\scriptsize [1156, 1315]} \\ {\scriptsize $\Delta = -14.2\%$}}
& \makecell{26.59 \\ {\scriptsize [22.2, 31.4]} \\ {\scriptsize $\Delta = -38.3\%$}}
& \makecell{451.5 \\ {\scriptsize [265.6, 669.9]} \\ {\scriptsize $\Delta = -83.0\%$}} \\
DH-60 ($\beta^{*} = 1$)
& \makecell{961 \\ {\scriptsize [934, 989]} \\ {\scriptsize $\Delta = -13.7\%$}}
& \makecell{\textbf{1191} \\ {\scriptsize [1150, 1234]} \\ {\scriptsize $\Delta = -16.8\%$}}
& \makecell{\textbf{9.39} \\ {\scriptsize [8.63, 10.23]} \\ {\scriptsize $\Delta = -78.2\%$}}
& \makecell{\textbf{2.86} \\ {\scriptsize [0.51, 5.86]} \\ {\scriptsize $\Delta = -99.9\%$}} \\
\bottomrule
\end{tabular}
\end{table*}

Table~\ref{tab:results_compact} shows that the operating point selected
on the simplified calibration network remains effective on the full
Shipai domain. Both dynamic strategies reduce tail evacuation time
relative to Static routing: $\mu_{50}[T_{0.95}]$ falls from \SI{1114}{s} to
\SI{956}{s} under Density-60 and \SI{961}{s} under DH-60, while
$\mu_{50}[T_{0.99}]$ falls from \SI{1431}{s} to \SI{1228}{s} and
\SI{1191}{s}, respectively. Direct paired analysis of the $50$ matched
scenarios gives $\mu_{50}[T_{0.95}^{\mathrm{DH}}-T_{0.95}^{\mathrm{Density}}]
=\SI{4.6}{s}$ (95\% paired-bootstrap CI:
\SI{-15.3}{s} to \SI{22.6}{s}) and
$\mu_{50}[T_{0.99}^{\mathrm{DH}}-T_{0.99}^{\mathrm{Density}}]
=\SI{-36.8}{s}$
(95\% CI: \SI{-122.1}{s} to \SI{31.3}{s}). Both intervals include
zero, so the data do not show evidence of a difference in mean run-level
tail evacuation time at the selected operating point. This result is not
an equivalence claim.

The safety outcomes separate much more strongly. Density-60 reduces
$\mu_{50}[\Sigma_s]$ by $38.3\%$ and $\mu_{50}[H]$ by $83.0\%$ relative to Static, showing
that congestion redistribution alone can incidentally reduce exposure.
However, DH-60 reduces $\mu_{50}[\Sigma_s]$ by $78.2\%$ and nearly eliminates
Hard Harm ($99.9\%$ reduction). Thus the Shipai benchmark does not support a
single ``faster means safer'' interpretation: the two dynamic strategies
are similar in clearance efficiency but qualitatively different in
exposure control. The same 50 paired Shipai runs are stratified by
hazard-source class in \appref{sm:disaster_locations}.


\subsection{Spatial redistribution mechanisms}
\label{sec:shipai_spatial_mechanisms}


\begin{figure}[!htbp]
  \centering
  \includegraphics[width=0.92\textwidth]{figures/07_shipai_case/fig7_spatial_reduction.pdf}
  \caption{Quantitative Static-relative reductions in the Shipai spatial
  fields. For each strategy, the fields are first averaged cellwise over
  the 50 paired Shipai scenarios, and all summaries are then restricted to
  the shared walkable-road mask $\mathcal{R}$. Panel (a) summarises the
  cumulative density-load field $L_{\rho}$ through
  $\sum_{\mathbf{p}\in\mathcal{R}}L_{\rho}(\mathbf{p})$,
  $Q^{+}_{0.95}(L_{\rho})$,
  and $N^{s}_{>0.99}(L_{\rho})$. Panel (b) summarises the
  congestion-duration field $T_{\mathrm{cong}}$ through
  $Q^{+}_{0.95}(T_{\mathrm{cong}})$ and
  $N^{s}_{>0.99}(T_{\mathrm{cong}})$.}
  \label{fig:shipai_spatial_reduction}
\end{figure}

To connect the evacuation outcomes with spatial mechanisms, we first
quantify the accumulated spatial fields over the 50 paired Shipai
scenarios. Figure~\ref{fig:shipai_spatial_reduction} summarises the
Static-relative reductions using the definitions in
\appref{sm:spatial_diagnostics}. Relative to
Static, the total cumulative density load
$\sum_{\mathbf{p}\in\mathcal{R}}L_{\rho}(\mathbf{p})$ falls by $3.4\%$
under Density-60 and $3.9\%$ under DH-60;
$Q^{+}_{0.95}(L_{\rho})$ falls by $14.2\%$ and $15.6\%$, respectively.
The same-field Static-threshold count $N^{s}_{>0.99}(L_{\rho})$ decreases
by $60.2\%$ for Density-60 and $63.8\%$ for DH-60. The
congestion-duration field shows the same pattern:
$Q^{+}_{0.95}(T_{\mathrm{cong}})$ decreases by $12.7\%$ and $15.1\%$,
while $N^{s}_{>0.99}(T_{\mathrm{cong}})$ contracts by $68.9\%$ and
$69.5\%$, respectively.

Figure~\ref{fig:shipai_maps} then provides a detailed spatial illustration
from a representative paired Shipai scenario. The same scenario is plotted
under Static, Density-60, and DH-60, with rows showing the route-use field
$F_{\mathrm{route}}(\mathbf{p})$, cumulative density-load field
$L_{\rho}(\mathbf{p})$, and congestion-duration field
$T_{\mathrm{cong}}(\mathbf{p})$. This figure is an illustrative
single-scenario map rather than the
source of the aggregate percentages in
Figure~\ref{fig:shipai_spatial_reduction}. In this example, all strategies
concentrate load on the main east--west corridors, boundary-connected
streets, and exit approaches, but the dynamic strategies visibly reduce the
intensity of the dominant Static hot-spots.

This spatial evidence clarifies why Density-60 and DH-60 have comparable
tail evacuation times. At the selected update period, both strategies
are similarly effective at dispersing persistent congestion away from
the most overloaded Static corridors. Their difference lies instead in
where the redistributed flow is allowed to go. Density-60 reacts only to
crowd state and can still assign attractive alternatives through
hazard-affected corridors. DH-60 raises the navigation cost of hazardous
cells and blocks cells beyond the hard threshold, preserving the
congestion-relief benefit while preventing exposure from accumulating
along unsafe routes. The strong reduction in $\mu_{50}[\Sigma_s]$ and
$\mu_{50}[H]$ is
therefore a hazard-selection effect layered on top of a similar
congestion-redistribution effect.

\appref{sm:disaster_locations} further stratifies the 50 paired Shipai
scenarios by hazard-source location: near exits (NE, $K_{\mathrm{NE}}=18$),
near main roads (NR, $K_{\mathrm{NR}}=18$), and far from both exits and
main roads (FR, $K_{\mathrm{FR}}=14$).
Writing $K_g$ for the sample size of location class $g$, the corresponding
summaries are $\mu_{K_g}[\cdot]$.
The location analysis shows that source placement has only a modest
effect on tail clearance but a large effect on exposure. Under Static
routing, the groupwise mean $\mu_{K_g}[T_{0.95}]$ remains between
\SI{1107}{s} and \SI{1126}{s} across the three classes, whereas
$\mu_{K_g}[\Sigma_s]$ ranges from
$33.53\times10^6$ to $60.59\times10^6$
risk-score--person--time-step units from NR to FR.
Hard Harm follows the same ordering, rising from
$1.5075\times10^6$ agent--time-step pairs in NR to
$4.2973\times10^6$ in FR. A hazard source
far from both exits and main roads is therefore not the slowest case in
clearance time, but it is the most severe case in accumulated risk.



\begin{figure}[!htbp]
  \centering
  \includegraphics[width=0.84\textwidth]{figures/07_shipai_case/compare_frequency_density_time_3x3_ppt_2.pdf}
  \caption{Representative spatial evacuation fields for one paired Shipai
  scenario. Columns compare Static, Density-60, and DH-60, respectively.
  Rows show three accumulated spatial diagnostics:
  $F_{\mathrm{route}}(\mathbf{p})$, $L_{\rho}(\mathbf{p})$, and
  $T_{\mathrm{cong}}(\mathbf{p})$.
  Each row preserves its own colour scale, enabling within-metric comparison
  of how the three routing strategies redistribute persistent corridor use,
  crowd load, and congestion duration within this illustrative simulation.
  Aggregate Static-relative reductions are quantified in
  Figure~\ref{fig:shipai_spatial_reduction}.}
  \label{fig:shipai_maps}
\end{figure}
\FloatBarrier

\subsection{Hazard-source location stratification}
\label{sec:shipai_source_location}


The same stratification confirms that the advantage of DH-60 is not
driven by a single favourable source class. Density-60 reduces
$\mu_{K_g}[\Sigma_s]$ in all three classes, but leaves substantial residual
exposure, especially for FR
($32.43\times10^6$ risk-score--person--time-step units). DH-60 reduces $\mu_{K_g}[\Sigma_s]$ by
$74.6$--$83.7\%$ across the three classes and nearly eliminates
Hard Harm, reaching a zero mean tally in NR and only
$4.5\times10^3$ and $4.4\times10^3$ agent--time-step pairs in NE and FR. The location
analysis therefore reinforces the main Shipai result: congestion-aware
rerouting is an efficiency intervention with partial safety benefits,
whereas density--hazard coupled rerouting is the risk-control mechanism
needed to suppress location-dependent exposure tails.


\subsection{Case-study synthesis}
\label{sec:shipai_synthesis}

DH-60 retains the congestion-relief effect of Density-60 while reducing
both Soft Risk and Hard Harm. The Shipai results therefore support treating
Density-60 as an efficiency baseline and DH-60 as the hazard-aware
alternative for controlling the model-based exposure indicators.


\FloatBarrier
\section{Discussion}
\label{sec:discussion}

The experiments in
Sections~\ref{sec:strategy_calibration}--\ref{sec:shipai} lead to three
main interpretations. First, congestion management and hazard-aware
routing address different parts of evacuation risk: the former improves
clearance, whereas the latter is needed to control exposure. Second, the
operating point calibrated on the simplified network remains effective
on the full Shipai benchmark, but its transfer should be read as
empirical support within the tested morphology rather than as a
universal parameter law. Third, some comparisons are inherently
scenario-dependent, especially the relative safety of Static and
Density-aware routing when hazard and congestion fields are arranged
differently in space.

\subsection{Efficiency--Safety Decoupling}
\label{sec:disc_coupling}

On the Shipai benchmark, the paired analysis does not detect a mean
tail-time difference between Density-60 and DH-60: the estimated
DH-60 minus Density-60 differences are
$\mu_{50}[T_{0.95}^{\mathrm{DH}}-T_{0.95}^{\mathrm{Density}}]
=\SI{4.6}{s}$
(95\% CI: \SI{-15.3}{s} to \SI{22.6}{s}) and
$\mu_{50}[T_{0.99}^{\mathrm{DH}}-T_{0.99}^{\mathrm{Density}}]
=\SI{-36.8}{s}$ (95\% CI: \SI{-122.1}{s} to \SI{31.3}{s}). These intervals
include zero but do not establish equivalence. The safety outcomes,
however, differ sharply. Relative to
Static, DH-60 reduces $\mu_{50}[\Sigma_s]$ by $78.2\%$ and
$\mu_{50}[H]$ by
$99.9\%$, compared with $38.3\%$ and $83.0\%$ under Density-60. Thus,
similar clearance performance can coexist with a large separation in
Soft Risk and Hard Harm. This is the main empirical reason for treating
efficiency and safety as distinct objectives rather than as substitutes.

The mechanism is the interaction between crowd redistribution and the
hazard field. A density-aware potential reacts to crowd state, so it can
move pedestrians away from congested corridors but has no direct
information about whether the alternative corridors are safer. When
congestion and hazard are spatially correlated, this can redirect flow
toward hazardous cells. The density--hazard potential reduces this risk
by raising the cost of high-intensity cells and excluding cells that
cross the blocking threshold. It therefore changes the route-choice
landscape before exposure accumulates, rather than only responding to
queues after they form.

This result also clarifies the role of the evaluation metrics. Mean
clearance time would miss the separation between Density-60 and DH-60.
Tail-sensitive and risk-sensitive run-level indicators
($T_{0.99}^{(k)}$, $\Sigma_s^{(k)}$, and $H^{(k)}$),
with their explicitly labelled ensemble summaries, expose the distinction between fast evacuation and
safe evacuation. The present evidence is based on scenario ensembles and
bootstrap confidence intervals; a fully probabilistic reliability
assessment would require a larger uncertainty campaign for converged
exceedance and CVaR estimates.

\subsection{What Transfers Across Scales, and What Does Not}
\label{sec:disc_transfer}

The simplified calibration network and the Shipai benchmark differ by
about one order of magnitude in agent count, several-fold in domain
area, and substantially in topological complexity. The comparison does
not prove morphology-independent generality, but it does identify which
patterns remain stable under the tested scale change.

The selected $(\beta^{*}, \Delta t_u^{*})$ retains the intended
behaviour on Shipai: the hazard weight lies near the knee of the
efficiency--exposure trade-off, the update period avoids the high
computational burden of \SI{10}{s} updates, and Density--Hazard yields the
lowest model-based exposure indicators among the tested alternatives. This supports the
two-stage workflow used in this study: parameter search can be carried
out on a tractable network, followed by verification on the target
settlement. The result should not be interpreted as a guarantee that the
same parameters are optimal for all urban forms; rather, it shows that
the calibrated balance is not destroyed by moving from the simplified
network to the Shipai layout.

On the simplified network, DH has a clear tail-time penalty relative to
Density-aware in $Q_{0.95,49}[T_{0.95}]$ ($+7.7\%$). On Shipai, the corresponding
point estimate is $+0.5\%$, and the paired confidence interval for the
absolute difference (\SI{-15.3}{s} to \SI{22.6}{s}) includes zero. One
plausible explanation is path redundancy. In a sparse network, hazard avoidance has few
available alternatives and therefore increases travel time. In the
Shipai layout, more competing path families can absorb hazard-driven
redirection at lower marginal time cost. The efficiency cost of
risk-aware routing is therefore not a fixed property of the method; it
depends on the redundancy and connectivity of the network.

The simplified network shows a counter-intuitive pattern:
Density-aware yields higher $Q_{0.95,49}[\Sigma_s]$ than Static
($2.10$ vs.\ $2.03 \times 10^{6}$). On Shipai, the ordering
reverses, with Density-60 reducing $\mu_{50}[\Sigma_s]$ by $38.3\%$ relative to
Static. This contrast is consistent with differences in the spatial
relationship between congestion and hazard. When density relief moves
pedestrians into hazard-affected corridors, congestion-aware routing can
increase exposure; when congestion and hazard are more weakly aligned,
density redistribution can also reduce some exposure incidentally. This
is why Density-aware guidance should not be treated as a safety
intervention unless exposure is evaluated directly.

At $\beta^{*}$, DH nearly eliminates Hard Harm on the simplified network
($\mu_{49}[H]=10.9$), and larger $\beta$ values eliminate it in that
calibration setting. On Shipai, a small residual remains
($\mu_{50}[H] \approx 2.86\times10^3$ agent--time-step pairs), although it is about three orders of magnitude lower
than under Static. This residual may partly reflect the finite update
period: within a \SI{60}{s} interval, transient crowd and hazard
conditions can evolve before the global field is refreshed. Reducing
$\Delta t_u$ could lower this residual further, but
Section~\ref{sec:update_period} shows that the computation cost rises
steeply below the selected operating point. The practical question is
therefore not whether risk can be reduced further, but how much residual
exposure is acceptable for a given computation budget.

\subsection{Limitations and Future Directions}
\label{sec:limitations}

The conclusions above support routing-level guidance under controlled
simulation conditions, but they rest on modelling choices that should be
relaxed in future work.

First, agents are assumed to comply fully with the navigation field;
behavioural heterogeneity, group interactions, and partial compliance
are not represented. The reported gains should therefore be read as the
performance of the guidance logic under ideal compliance. Combining the
framework with empirically calibrated route-choice, group-behaviour, and
partial-compliance models would allow the same routing policy to be tested
under more realistic behavioural response; the cited dynamic-guidance
studies are not used here as evidence of such compliance models.

Second, the hazard model is parameterised rather than physics-based.
Detailed processes---turbulence, smoke stratification, combustion
chemistry, multi-hazard coupling---are absorbed into the scalar
intensity $R$ and thresholds such as $R_H$ and
$R_{\mathrm{block}}$. The reported conclusions are
planning-level evidence on the \emph{relative} performance of routing
strategies under controlled hazard scenarios, not absolute predictions
of fire or toxic-exposure outcomes. Two-way coupling with physics-based
simulators would improve hazard realism, but would also increase the
computational burden of online guidance.
One-factor parameter sweeps illustrating the response of this simplified
hazard field are reported in \appref{sm:hazard_sensitivity}. They do not
evaluate the sensitivity of evacuation outcomes or strategy rankings;
such an outcome-level robustness analysis remains future work.

Third, the empirical reliability claim rests on scenario ensembles and
tail-sensitive statistics rather than on converged exceedance or CVaR
estimates. A larger Monte Carlo campaign with explicit uncertainty
propagation through the hazard, demand, and behavioural channels is
needed to support fully probabilistic reliability statements; the
framework is compatible with such an analysis but does not itself
establish it. The source-location stratification in
\appref{sm:disaster_locations} is therefore best read as a targeted
stress test rather than a converged uncertainty propagation study.

Beyond these three, the two-dimensional formulation does not represent
vertical evacuation or terrain effects, restricting application to flat
ground-level settlements; and while Shipai is a representative
high-density urban village, the cross-morphology generality of
$(\beta^{*}, \Delta t_u^{*})$---across high-rise complexes, transit
hubs, or open public squares---remains an open empirical question. The
simplified-vs.-Shipai comparison reported here suggests that the
safety advantage of hazard-aware routing can transfer while absolute
risk levels remain morphology- and hazard-placement-dependent. Broader
cross-scenario studies are needed to test this pattern directly.


\section{Conclusion}
\label{sec:conclusion}

This study developed a density--hazard coupled global navigation
framework for pedestrian evacuation in maze-like urban-village
networks. The framework links an empirically calibrated microscopic
movement model, an online global potential that combines congestion and
hazard costs, and a two-stage protocol that calibrates operating
parameters on a simplified network before verification on the full
Shipai benchmark.

The central result is that evacuation efficiency and model-based exposure outcomes are
not interchangeable objectives. On Shipai, the paired confidence
intervals for the Density-60--DH-60 tail-time differences include zero,
so the analysis does not detect a difference in mean tail evacuation
time; it does not establish equivalence. DH-60 nevertheless reduces
Soft Risk by $78.2\%$ and Hard Harm by $99.9\%$ relative to
Static, compared with $38.3\%$ and $83.0\%$ under Density-60. This
separation shows that congestion-aware routing can improve clearance
without adequately controlling hazard exposure; explicit hazard
awareness in the navigation potential is needed to suppress risk at the
calibrated operating point.

The calibration results identify a practical operating point:
$(\beta^{*}, \Delta t_u^{*}) = (1.0,\,\SI{60}{s})$. At this point,
the hazard weight sits near the knee of the efficiency--exposure
trade-off, and the update period substantially reduces field-update cost
relative to \SI{10}{s} updates while preserving most of the
evacuation benefit. Its successful application to Shipai---a larger,
more irregular network with substantially more agents---supports the use
of simplified-network calibration as a screening step before expensive
full-scale simulation, while leaving absolute risk levels dependent on
network morphology and hazard placement.

For evacuation planning in dense informal settlements, the implication
is direct: strategies should be judged with tail-sensitive time,
Soft Risk, and Hard Harm indicators
($T_{0.99}^{(k)}$, $\Sigma_s^{(k)}$, and $H^{(k)}$),
reported with an explicitly identified ensemble functional, not with mean
clearance time alone. The conclusions should be
interpreted within the assumptions of full compliance, parameterised
hazard intensity, and scenario-ensemble rather than fully converged
probabilistic reliability evidence. Future work should couple the
framework with compliance-aware behaviour, physics-based hazard
simulation, larger uncertainty campaigns, and cross-morphology tests.



\FloatBarrier   % 放在 \subsection 结束处，阻止后面的浮动跨过去
% ------------------------------------------------------------------------
\newpage  % 另起一页
\appendix  % 附录部分 （1）灾害演变 （2）人口初始化 （3）灾害发生位置


% &&&&&&&&&&&&&&&&&&&&&&&&&&&&&&&&&&&&&&&&&&&&&&&&&&&&&&&&&&&&&&&

\section{SPH Density-Estimator Sanity Check}
\label{app:sph_density}

This appendix documents the sanity check of the SPH-based density estimator
used in the main text. The checked object is the continuous Poly6 estimate
with the same support radius $h=1.2$\,m used by the simulation; it is not a
second raster kernel. The objective is to establish a quantitative mapping
between the estimated density $\hat{\rho}_{\mathrm{SPH}}$ and the physical
pedestrian density $\rho_{\mathrm{real}}$ (ped/m$^2$), thereby ensuring
physical interpretability and numerical consistency.


% ---------------------------------------------------------
\subsection{Problem Formulation}
\label{app:calib_problem}

In the SPH framework, local density is estimated as
\begin{equation}
\hat{\rho}_{\mathrm{SPH}}
=
\sum_j m_j \, W_{\mathrm{poly6}}(r_{ij}, h),
\end{equation}
where $W_{\mathrm{poly6}}$ is the Poly6 kernel and $h=1.2$\,m is the
smoothing length used in both this check and the simulation source code.

Although the kernel is normalised, the discrete nature of particle sampling,
finite domain effects, and local spatial variability may introduce deviations
from the physical density. Therefore, we fit the diagnostic relationship
\begin{equation}
\rho_{\mathrm{real}}
=
a_\rho \, \hat{\rho}_{\mathrm{SPH}} + b_\rho,
\label{eq:calib_relation_app}
\end{equation}
and evaluate whether $a_\rho \approx 1$ and $b_\rho \approx 0$ hold in practice.


% ---------------------------------------------------------
\subsection{Experimental Setup}
\label{app:calib_setup}

The check is conducted in a square domain of size
$L \times L = 10\,\mathrm{m} \times 10\,\mathrm{m}$.
The Poly6 support radius is fixed at $h=1.2$\,m, matching the simulation,
and is not fitted in this experiment.
The target density range is
\[
\rho_{\mathrm{real}} \in [0.1,\,4.5] \ \mathrm{ped/m}^2.
\]

For each density level, two spatial configurations are considered:

\begin{itemize}
    \item \textbf{Uniform placement:}
    pedestrians are arranged on a regular grid, representing an idealised
    homogeneous configuration.

    \item \textbf{Random hard-core placement:}
    pedestrians are sampled randomly subject to a minimum inter-agent distance
    $d_{\min}$, preventing unrealistic overlap and approximating real crowd
    distributions.
\end{itemize}

In this study, $d_{\min}=0.3$\,m is used unless otherwise stated.


% ---------------------------------------------------------
\subsection{Sampling Strategy}
\label{app:sampling_strategy}

To reduce spatial sampling bias, density is not evaluated at a single point.
Instead, we use a grid of sampling centres located within the interior domain.
Specifically, an $8\times8$ lattice of sampling points is constructed, all at
least one smoothing length ($1.2$\,m) away from the boundary.

The SPH density is then averaged over these centres:
\begin{equation}
\bar{\rho}_{\mathrm{SPH}}
=
\frac{1}{M}
\sum_{k=1}^{M}
\hat{\rho}_{\mathrm{SPH}}(\mathbf{p}_k),
\end{equation}
where $M=64$ is the number of sampling centres.

For the random hard-core configuration, $N_r=50$ independent realisations are
generated at each density level. The final SPH density is computed as
\begin{equation}
\langle \rho_{\mathrm{SPH}} \rangle
=
\frac{1}{N_r}
\sum_{s=1}^{N_r}
\bar{\rho}_{\mathrm{SPH}}^{(s)}.
\end{equation}

This two-level averaging (spatial + ensemble) reduces variance
caused by local clustering or voids.


% ---------------------------------------------------------
\subsection{Effective Density}
\label{app:effective_density}

Due to the hard-core constraint, the number of placed agents may be slightly
lower than the target value at high densities. Therefore, the effective density
is defined as
\begin{equation}
\rho_{\mathrm{real}}
=
\frac{N_{\mathrm{actual}}}{L^2},
\end{equation}
where $N_{\mathrm{actual}}$ is the actual number of successfully placed agents.
This ensures that the check is based on physically realizable configurations.


% ---------------------------------------------------------
\subsection{Results and Analysis}
\label{app:calib_results}

The check results show a strong linear relationship between SPH density
and physical density under both placement strategies.

For the random hard-core configuration, the fitted relation is
\begin{equation}
\rho_{\mathrm{real}}
=
1.020\,\hat{\rho}_{\mathrm{SPH}} - 0.019,
\qquad
R^2 = 1.000.
\end{equation}

The slope is close to unity and the intercept is small, indicating that
the SPH estimator provides an approximately unbiased estimate of the
physical density over the tested range.

For the uniform configuration, a similar linear trend is observed, though with
a slightly larger slope and small positive offset. This deviation is attributed
to the periodic structure of the grid, which introduces mild interaction
between particle spacing and kernel support.

Overall, the results confirm that the SPH density is a reliable proxy for
physical density across both ordered and disordered configurations. This
result validates the $h=1.2$\,m Poly6 estimator. The subsequent
$\Delta_\rho=1.0$\,m max aggregation used to construct the navigation
raster is a separate deterministic mapping and is not assigned an
additional smoothing parameter.


% ---------------------------------------------------------
\subsection{Discussion}
\label{app:calib_discussion}

The sanity check highlights several observations:

\begin{itemize}
    \item The SPH density exhibits a near-linear mapping to physical density
    over a wide operational range.

    \item Random hard-core configurations provide a more realistic reference
    than uniform grids, as they capture spatial disorder and non-overlapping
    constraints inherent in pedestrian systems.

    \item Multi-centre and multi-realisation averaging are essential to reduce
    variance and obtain stable diagnostic estimates.

    \item Residual deviations at high densities are primarily caused by packing
    constraints rather than kernel inaccuracies.
\end{itemize}

These findings justify the direct use of $\hat{\rho}_{\mathrm{SPH}}$ in the main
model without additional correction, as adopted in the main text.








% =========================================================
% Appendix B: Hazard Intensity Model -- Parameter-Response Illustration
% =========================================================
% =========================================================
\section{Hazard Intensity Model: Parameter-Response Illustration}
\label{sm:hazard_sensitivity}

This appendix documents one-factor parameter-response illustrations for
the hazard model in \S\ref{sec:hazard_field}. The aim is to show how
arrival time, intensity growth, and hard blocking jointly determine the
constructed hazard field. The sweeps display temporal response curves
and field snapshots only; they do not propagate each parameter setting
through evacuation simulations and therefore do not establish the
sensitivity or robustness of $T_{0.95}$, exposure, $H$, or strategy
rankings.


% % =========================================================
% \subsubsection{Hazard-Field Propagation via Arrival Time and Time-Response Intensity}

% We represent hazard spread on the raster grid, where each cell~$\mathbf{p}$ is
% assigned a medium-dependent propagation speed $v_h(\mathbf{p}) > 0$ (e.g.,
% different values for open space vs.\ buildings; buildings can optionally be
% treated as hard barriers). Given a set of ignition/source cells
% $\mathcal{S} \subset \mathcal{G}$, we first compute an \emph{arrival-time} field
% $T_h(\mathbf{p})$ and then derive a \emph{time-dependent intensity}
% $R(\mathbf{p}, t) \in [0, R_{\max}]$ and a \emph{hard-blocking mask}
% $B(\mathbf{p}, t) \in \{0, 1\}$.



% ---------------------------------------------------------
\subsection{Arrival-Time Field}
\label{sec:hazard_arrival_time}

The arrival time $T_h(\mathbf{p})$ encodes the earliest time the hazard front
can reach each cell and is defined as the minimum travel time over all paths
$\pi_h$ from any source cell:
\begin{equation}
T_h(\mathbf{p}) = \min_{\pi_h:\,\mathcal{S} \rightarrow \mathbf{p}}
\int_{\pi_h} \frac{1}{v_h(\boldsymbol{\xi})} \, \mathrm{d}\ell.
\label{eq:traveltime}
\end{equation}
In the continuum, this corresponds to the isotropic Eikonal equation
\begin{equation}
\|\nabla T_h(\mathbf{p})\| = \frac{1}{v_h(\mathbf{p})},
\qquad T_h(\mathbf{p}) = 0, \; \forall \mathbf{p} \in \mathcal{S},
\label{eq:eikonal}
\end{equation}
whose viscosity solution yields the first-arrival time of a monotonically
advancing front. This formulation and its fast solvers are standard in
front-propagation and level-set literature
\citep{Sethian1996FastMarching,Zhao2005FastSweeping}.

On the raster, we discretise Eq.~\eqref{eq:eikonal} using an 8-neighbourhood
$\mathcal{N}_8(\mathbf{p})$ and solve via iterative Bellman-type relaxation:
\begin{equation}
T_h^{(k+1)}(\mathbf{p}) \leftarrow \min\!\Bigg(
T_h^{(k)}(\mathbf{p}), \;
\min_{\mathbf{q} \in \mathcal{N}_8(\mathbf{p})}
\bigg[ T_h^{(k)}(\mathbf{q}) + \frac{d(\mathbf{p}, \mathbf{q})}{v_h(\mathbf{p})} \bigg]
\Bigg),
\label{eq:relax_update}
\end{equation}
where $d(\mathbf{p}, \mathbf{q}) \in \{\Delta, \sqrt{2}\Delta\}$ is the grid
distance. This iterative perspective is closely related to sweeping/relaxation-type
Eikonal solvers \citep{Zhao2005FastSweeping}. In our implementation, $v_h(\mathbf{p})$ is
medium-dependent (e.g., $v_{\mathrm{open}}$ vs.\ $v_{\mathrm{build}}$), consistent
with classical spread-rate modelling \citep{Rothermel1972}.
The arrival-time map is computed \emph{once} per hazard configuration and cached;
it is recomputed if either the source locations or the propagation-speed
field $v_h(\mathbf{p})$ changes.


% ---------------------------------------------------------
\subsection{Time-Dependent Hazard Intensity and Hard Blocking}
\label{sec:hazard_intensity}

Given $T_h(\mathbf{p})$, we define a time-dependent intensity
$R(\mathbf{p}, t) \in [0, R_{\max}]$ as an arrival-time-referenced response.
Let $s = t - T_h(\mathbf{p})$ denote the relative time, where $s < 0$ indicates
\emph{pre-arrival} and $s \geq 0$ indicates \emph{post-arrival}. We model a
piecewise exponential response:
\begin{equation}
R(\mathbf{p}, t) =
\begin{cases}
R_0 \,\exp\!\left(\dfrac{s}{\tau_{\mathrm{pre}}}\right),
& s < 0, \\[8pt]
R_0 + (R_{\max} - R_0)\!\left(1 - \exp\!\left(-\dfrac{s}{\tau_{\mathrm{rise}}}\right)\right),
& s \geq 0,
\end{cases}
\label{eq:intensity_piecewise}
\end{equation}
where $R_0 = \alpha_{\mathrm{pre}} R_{\max}$ with $\alpha_{\mathrm{pre}} \in [0,1]$
controls the intensity at the moment of arrival, $\tau_{\mathrm{pre}}$ governs the
pre-arrival exposure decay, and $\tau_{\mathrm{rise}}$ governs the post-arrival
rise timescale. The pre-arrival branch models early leakage ahead of the hazard
front (e.g., smoke diffusing before the flame arrives), while the post-arrival
branch models gradual intensification toward~$R_{\max}$.
Such time-evolving hazard fields are consistent with engineering fire modelling
frameworks \citep{McGrattan2013FDS}, and dose-accumulation models such as the
Fractional Effective Dose (FED) provide a principled basis for connecting
concentration--time histories to incapacitation thresholds \citep{Purser1989FED}.
In the present model, $R$ is not derived from gas concentration, heat, or
FED; these quantities therefore represent a future physical-calibration
pathway rather than a current conversion from $R$ to health outcomes.
One-factor response curves for $\tau_{\mathrm{rise}}$,
$\tau_{\mathrm{pre}}$, and $\alpha_{\mathrm{pre}}$ are provided below.

We further define a \emph{hard-blocking mask} $B(\mathbf{p}, t) \in \{0, 1\}$
that disables traversal once the hazard exceeds a blocking threshold
$R_{\mathrm{block}}$:
\begin{equation}
B(\mathbf{p}, t)
= \mathbb{I}\!\Big(
t \geq T_h(\mathbf{p}) \;\wedge\; R(\mathbf{p}, t) \geq R_{\mathrm{block}}
\Big).
\label{eq:blocking}
\end{equation}
From Eq.~\eqref{eq:intensity_piecewise}, the blocking onset time admits a closed
form:
\begin{equation}
t_{\mathrm{block}}(\mathbf{p})
= T_h(\mathbf{p}) + \tau_{\mathrm{rise}} \ln\!\left(
\frac{R_{\max} - R_0}{R_{\max} - R_{\mathrm{block}}}
\right),
\label{eq:block_time}
\end{equation}
which clarifies how $\tau_{\mathrm{rise}}$, $\alpha_{\mathrm{pre}}$, and
$R_{\mathrm{block}}$ jointly control the onset of impassable regions.
Equation~\eqref{eq:block_time} applies when
$R_0<R_{\mathrm{block}}<R_{\max}$. If
$R_{\mathrm{block}}\le R_0$, blocking begins at front arrival
($t_{\mathrm{block}}=T_h$); if
$R_{\mathrm{block}}\ge R_{\max}$, no finite post-arrival blocking time
exists. In all reported experiments, $R_0=3.0$,
$R_{\mathrm{block}}=6.0$, $R_{\max}=10$, and
$\tau_{\mathrm{rise}}=120$\,s, giving
$t_{\mathrm{block}}-T_h\approx67.2$\,s.

The extended intensity response in Eq.~\eqref{eq:intensity_piecewise} is
governed by four parameter groups: the post-arrival ramp-up
$\tau_{\mathrm{rise}}$, the pre-arrival leakage $\tau_{\mathrm{pre}}$, the
arrival-level fraction $\alpha_{\mathrm{pre}}$, and the medium-dependent
propagation speeds $(v_{\mathrm{open}}, v_{\mathrm{build}})$. This appendix
characterises each parameter's role through (i)~one-dimensional temporal
response curves and (ii)~a representative spatial footprint comparison, and
documents the 20 preset cases generated for field-level illustration.


% ---------------------------------------------------------
\subsection{Baseline Configuration}
\label{sm:baseline_config}

All one-factor sweeps are conducted around a common baseline
(\texttt{custom2}: $v_{\mathrm{open}} = 0.10$\,m/s,
$v_{\mathrm{build}} = 0.05$\,m/s, $R_{\max} = 10$,
$R_{\mathrm{block}} = 6.0$,
$\tau_{\mathrm{pre}} = 500$\,s, $\tau_{\mathrm{rise}} = 120$\,s,
$\alpha_{\mathrm{pre}} = 0.30$, buildings non-blocking).
This configuration is chosen because it activates both pre-arrival leakage and
post-arrival growth in a non-degenerate regime (neither step-like nor nearly
flat), ensuring identifiability of each parameter's effect in one-factor sweeps.
The full 20-case parameter matrix is listed in Table~\ref{tab:sm_presets}.


% ---------------------------------------------------------
\subsection{Temporal Response Curves}
\label{sm:temporal_curves}

Figures~\ref{fig:sm_tau_rise}--\ref{fig:sm_speed} show the intensity
$R(\mathbf{x}, t)$ as a function of relative time $s = t - T_h(\mathbf{x})$
under one-factor sweeps. The arrival event is anchored at $s = 0$.

\begin{figure}[htbp]
    \centering
    \includegraphics[width=0.80\textwidth]{figures/appendix_b/intensity_sweep_tau_rise.png}
    \caption{Effect of $\tau_{\mathrm{rise}}$ on post-arrival intensity ramp-up.
    Smaller $\tau_{\mathrm{rise}}$ produces faster escalation (shorter grace
    period after $t = T_h$); larger values yield smoother saturation.}
    \label{fig:sm_tau_rise}
\end{figure}

\begin{figure}[htbp]
    \centering
    \includegraphics[width=0.80\textwidth]{figures/appendix_b/intensity_sweep_tau_pre.png}
    \caption{Effect of $\tau_{\mathrm{pre}}$ on pre-arrival leakage.
    Larger $\tau_{\mathrm{pre}}$ implies slower decay ahead of the front,
    producing stronger early exposure and reducing the effective safe-time
    margin before nominal arrival.}
    \label{fig:sm_tau_pre}
\end{figure}

\begin{figure}[htbp]
    \centering
    \includegraphics[width=0.80\textwidth]{figures/appendix_b/intensity_sweep_alpha_pre.png}
    \caption{Effect of $\alpha_{\mathrm{pre}}$ on arrival-level intensity
    ($R_0 = \alpha_{\mathrm{pre}} R_{\max}$). Higher $\alpha_{\mathrm{pre}}$
    lifts the entire post-arrival trajectory and increases the
    time-integrated intensity at a fixed observation point and time window.}
    \label{fig:sm_alpha_pre}
\end{figure}

\begin{figure}[htbp]
    \centering
    \includegraphics[width=0.80\textwidth]{figures/appendix_b/intensity_sweep_speed_open.png}
    \caption{Effect of propagation speed $v_{\mathrm{open}}$ on the temporal
    response at a fixed observation point. Faster propagation shifts the
    arrival time $T_h$ earlier (curve shifts left), shrinking the available
    evacuation window. The temporal shape of the response
    ($\tau_{\mathrm{pre}}$, $\tau_{\mathrm{rise}}$) remains unchanged; only
    the response is shifted earlier without changing its shape in relative time.}
    \label{fig:sm_speed}
\end{figure}

The four parameter groups separate hazard-field variation into distinct,
interpretable dimensions:
\begin{enumerate}
    \item $v_{\mathrm{open}},\, v_{\mathrm{build}}$: control \emph{when} the
    hazard reaches each location (arrival-time geometry $T_h$).
    \item $\alpha_{\mathrm{pre}}$: controls \emph{how severe} conditions are at
    the moment of arrival.
    \item $\tau_{\mathrm{rise}}$: controls \emph{how quickly} conditions
    escalate after arrival.
    \item $\tau_{\mathrm{pre}}$: controls \emph{how early} risk emerges before
    the front arrives.
\end{enumerate}
This factorisation enables a compact yet representative scenario suite: a small
number of ignition locations combined with 2--3 speed levels and a few stress
settings of $\{\alpha_{\mathrm{pre}}, \tau_{\mathrm{rise}},
\tau_{\mathrm{pre}}\}$ span both ``sudden-escalation'' and
``early-leakage'' field regimes. These presets define candidate hazard
fields; no claim of evacuation-outcome robustness is made from these
illustrations alone.


% ---------------------------------------------------------
\subsection{Spatial Footprint: Propagation Speed Sweep}
\label{sm:spatial_footprint}

Among the four parameter groups, propagation speed is the only sweep that
changes the arrival-time surface $T_h(\mathbf{x})$. The other three sweeps
($\tau_{\mathrm{rise}}$, $\tau_{\mathrm{pre}}$, $\alpha_{\mathrm{pre}}$) leave
$T_h$ unchanged, but they do not leave the full spatial hazard field
unchanged. At a fixed global time $t^*$, the relative time
$s=t^*-T_h(\mathbf{x})$ varies across space; changing any of these response
parameters therefore changes the contours of $R(\mathbf{x},t^*)$ and can
change the area satisfying $R\ge R_{\mathrm{block}}$. Figure~\ref{fig:sm_groupS_spatial}
shows Group~S as a representative spatial illustration because it additionally
changes arrival-time geometry. Spatial panels for Groups~A--C are not shown,
and the one-dimensional curves should not be interpreted as evidence that
their spatial intensity and blocking footprints are invariant.

\begin{figure}[htbp]
    \centering
    \includegraphics[width=1.0\linewidth]{figures/appendix_b/disaster_groupS_arrival_intensity_hardmask_stack3x5.pdf}
    \caption{Group~S (propagation speed sweep): field-level response of the
    hazard footprint. Each column corresponds to a different speed level
    (S01--S05, increasing $v_{\mathrm{open}}$ from left to right). Rows show
    (top)~arrival-time field $T_h(\mathbf{x})$, (middle)~intensity
    $R(\mathbf{x}, t^*)$ at the common evaluation time
    $t^*=\SI{1800}{s}$, and (bottom)~the hard mask at
    $R_{\mathrm{block}}=6.0$. Faster propagation
    compresses arrival times and expands the near-saturated hazardous region at
    the same evaluation instant.}
    \label{fig:sm_groupS_spatial}
\end{figure}

% =========================================================
% Appendix B.6: Preset Parameter Cases (alternative layout -- guaranteed no overflow)
% =========================================================

\subsection{Preset Parameter Cases}
\label{sm:preset_cases}

Table~\ref{tab:sm_presets} lists the 20 cases organised into four one-factor
sweep groups. Each group varies a single parameter (shown in \textbf{bold})
from the baseline while holding all others fixed.

\begin{table}[htbp]
\centering
\small
\renewcommand{\arraystretch}{1.20}
\caption{Twenty preset hazard-parameter cases for one-factor field-response
illustration; evacuation outcomes and strategy rankings are not evaluated in
these sweeps.
Baseline (\texttt{custom2}):
$v_{\mathrm{open}}\!=\!0.10$\,m/s,
$v_{\mathrm{build}}\!=\!0.05$\,m/s,
block\,$=\!0$,
$R_{\max}\!=\!10$,
$R_{\mathrm{block}}\!=\!6.0$,
$\tau_{\mathrm{pre}}\!=\!500$\,s,
$\tau_{\mathrm{rise}}\!=\!120$\,s,
$\alpha_{\mathrm{pre}}\!=\!0.30$.
All unvaried parameters inherit these baseline values.}
\label{tab:sm_presets}

\begin{tabularx}{\textwidth}{l l X}
\toprule
\textbf{Cases} & \textbf{Swept parameter} & \textbf{Values tested} \\
\midrule
A01--A05
  & $\tau_{\mathrm{rise}}$ (post-arrival ramp-up)
  & \textbf{20, 80, 120, 180, 300}\,s \\[4pt]
B01--B05
  & $\tau_{\mathrm{pre}}$ (pre-arrival leakage)
  & \textbf{100, 300, 500, 700, 900}\,s \\[4pt]
C01--C05
  & $\alpha_{\mathrm{pre}}$ (arrival intensity fraction)
  & \textbf{0.10, 0.30, 0.50, 0.70, 0.90} \\[4pt]
S01--S05
  & $(v_{\mathrm{open}},\, v_{\mathrm{build}})$ (propagation speed)
  & $v_{\mathrm{open}}=\textbf{0.01, 0.05, 0.10, 0.50, 1.00}$\,m/s;
    $v_{\mathrm{build}}=v_{\mathrm{open}}/2$ \tabularnewline
\bottomrule
\end{tabularx}

\vspace{4pt}
\footnotesize\textit{Note.}
Speed sweep (Group~S) uses direct propagation-speed values while
holding $v_{\mathrm{build}} = v_{\mathrm{open}}/2$. For example, S01 uses
$v_{\mathrm{open}} = 0.01$\,m/s, $v_{\mathrm{build}} = 0.005$\,m/s.
Fixed parameters for all groups:
block\,$= 0$, $R_{\max} = 10$, and $R_{\mathrm{block}}=6.0$.
\end{table}



\clearpage
% =========================================================
% Appendix C: Disaster Location Analysis
% =========================================================
\section{Disaster Location Analysis: Tail-Risk Dispersion}
\label{sm:disaster_locations}

This appendix examines whether the Shipai results depend on where the
hazard source is initialised. The analysis uses the same 50 random scenario
blocks for the three Shipai strategies in the main text: Static, Density-60,
and the density--hazard coupled strategy DH-60
($\beta=1$, $\Delta t_u=\SI{60}{s}$). Each run index shares the same demand
realisation and hazard source across the three strategies, allowing a
controlled comparison by source location.

\subsection{Experimental Design}

The random-input table classifies each Shipai hazard source into three
location types: near an exit (NE), near a main road (NR), or far from both
exits and main roads (FR). The resulting sample sizes are
$K_{\mathrm{NE}}=K_{\mathrm{NR}}=18$ and $K_{\mathrm{FR}}=14$ runs per
strategy. For each run we extract $T_{0.95}$, $T_{0.99}$, Soft Risk
$\Sigma_s$, and Hard Harm $H$. Mean values and uncertainty
bands are computed within each strategy--location cell using bootstrap
resampling. For a location class $g$ containing $K_g$ runs, all table
entries use the groupwise cross-run mean
$\mu_g[\cdot]\equiv\mu_{K_g}[\cdot]$.

\begin{table}[htbp]
\centering
\footnotesize
\setlength{\tabcolsep}{3.5pt}
\renewcommand{\arraystretch}{1.12}
\caption{Shipai performance metrics stratified by hazard-source location.
Every metric is summarised by $\mu_g[\cdot]$ within the indicated
strategy--location cell. $\Sigma_s$ is reported in $10^6$
risk-score--person--time-step units and
$H$ in $10^3$ agent--time-step pairs.}
\label{tab:sm_hazard_location}
\begin{tabular}{llrrrrr}
\toprule
\textbf{Location} & \textbf{Strategy} & \textbf{$K_g$}
  & \textbf{$\mu_g[T_{0.95}]$} & \textbf{$\mu_g[T_{0.99}]$}
  & \textbf{$\mu_g[\Sigma_s]$} & \textbf{$\mu_g[H]$ ($\times10^3$)} \\
\midrule
NE & Static      & 18 & 1107 & 1425 & 39.09 & 2509.9 \\
NE & Density-60 & 18 &  969 & 1310 & 27.22 &  604.0 \\
NE & DH-60       & 18 &  963 & 1193 &  9.88 &    4.5 \\
\midrule
NR & Static      & 18 & 1126 & 1456 & 33.53 & 1507.5 \\
NR & Density-60 & 18 &  955 & 1165 & 21.42 &  232.6 \\
NR & DH-60       & 18 &  965 & 1215 &  8.53 &    0.0 \\
\midrule
FR & Static      & 14 & 1107 & 1409 & 60.59 & 4297.3 \\
FR & Density-60 & 14 &  942 & 1201 & 32.43 &  536.7 \\
FR & DH-60       & 14 &  953 & 1157 &  9.86 &    4.4 \\
\bottomrule
\end{tabular}
\end{table}

\subsection{Location Effects on Tail Clearance}

Figure~\ref{fig:sm_hz_bars} reports the four primary metrics by source
location, and Figure~\ref{fig:sm_hz_interaction} shows the corresponding
strategy--location interaction profiles. Tail clearance is only weakly
moderated by hazard location. Under Static routing, $\mu_g[T_{0.95}]$ remains
between 1107 and 1126\,s across all three source classes, a spread of only
19\,s. This stability indicates that the 95th-percentile clearance time is
controlled mainly by network-wide demand and exit bottlenecks rather than
by the local position of the hazard source.

Dynamic guidance improves tail clearance in all three source classes.
Density-60 reduces $\mu_g[T_{0.95}]$ by 12.5\,\% in NE, 15.2\,\% in NR, and
14.9\,\% in FR. DH-60 gives a similar range of $\mu_g[T_{0.95}]$ improvements
(13.0--14.3\,\%), indicating that DH-60 retains similar Static-relative
clearance gains at the selected operating point. The
$\mu_g[T_{0.99}]$ patterns are slightly more location-sensitive: Density-60 is
strongest in NR (19.9\,\% reduction), whereas DH-60 is strongest in FR
(17.9\,\% reduction). Even so, the dominant result is that both dynamic
strategies consistently reduce tail clearance relative to Static routing.

\begin{figure*}[htbp]
  \centering
  \includegraphics[width=\textwidth]{figures/appendix_c/fig_hz1_grouped_bars.pdf}
  \caption{Evacuation and risk metrics stratified by hazard-source
  location. NE = near exit; NR = near main road; FR = far from exit and
  road. Error bars denote 95\,\% bootstrap confidence intervals.
  \textbf{(a)} $\mu_g[T_{0.95}]$; \textbf{(b)} $\mu_g[T_{0.99}]$;
  \textbf{(c)} $\mu_g[\Sigma_s]$; \textbf{(d)} $\mu_g[H]$ in
  unscaled agent--time-step pairs.}
  \label{fig:sm_hz_bars}
\end{figure*}

\begin{figure*}[htbp]
  \centering
  \includegraphics[width=\textwidth]{figures/appendix_c/fig_hz2_interaction.pdf}
  \caption{Strategy--location interaction plots for tail clearance and
  exposure metrics. Shaded bands denote 95\,\% bootstrap confidence
  intervals. Tail clearance varies modestly by source class, whereas the
  risk metrics respond strongly to hazard-source placement.}
  \label{fig:sm_hz_interaction}
\end{figure*}

\subsection{Location Effects on Soft Risk and Hard Harm}

Risk outcomes are much more sensitive to source location than clearance
times. Under Static routing, FR is the most severe class:
$\mu_g[\Sigma_s]$ reaches $60.59\times10^6$ risk-score--person--time-step
units, compared with $39.09\times10^6$ for NE and $33.53\times10^6$ for NR.
The Hard Harm tally follows the same ordering,
increasing from $1.5075\times10^6$ agent--time-step pairs in NR to
$4.2973\times10^6$ in FR. In other words, a
source that is remote from both exits and main roads produces the largest
tail exposure, even though it does not produce the longest $\mu_g[T_{0.95}]$.
This decoupling is the central diagnostic value of the location analysis:
clearance time alone does not identify the most hazardous source geometry.

Density-60 reduces exposure in every location class, but the reductions are
partial. Relative to Static routing, $\mu_g[\Sigma_s]$ falls by 30.4\,\% in NE,
36.1\,\% in NR, and 46.5\,\% in FR; Hard Harm falls by 75.9\,\%, 84.6\,\%,
and 87.5\,\%, respectively. These gains confirm that density redistribution
can incidentally lower exposure, especially when congested routes overlap
with hazardous regions. However, the residual FR exposure remains
$32.43\times10^6$ risk-score--person--time-step units, showing that congestion feedback alone cannot
reliably detect and avoid an isolated hazard pocket.

DH-60 gives a qualitatively different risk profile. It reduces $\mu_g[\Sigma_s]$
by 74.7\,\% in NE, 74.6\,\% in NR, and 83.7\,\% in FR
(Figure~\ref{fig:sm_hz_reduction}). The FR case is especially informative:
Soft Risk falls from $60.59\times10^6$ risk-score--person--time-step units
under Static routing to $9.86\times10^6$ under DH-60, a 6.1-fold reduction relative to
Static and a 3.3-fold reduction relative to Density-60. Hard Harm
is nearly eliminated in all three source classes:
DH-60 reaches a zero mean tally in NR and leaves only
$4.5\times10^3$ and $4.4\times10^3$ agent--time-step pairs in NE and FR.

\begin{figure*}[htbp]
  \centering
  \includegraphics[width=\textwidth]{figures/appendix_c/fig_hz3_reduction.pdf}
  \caption{Percentage improvement relative to the Static baseline by
  strategy and hazard-source location. Positive values indicate lower tail
  clearance time, Soft Risk, or Hard Harm. The DH-60 panel shows the largest
  Soft Risk reduction in FR and complete mean Hard Harm elimination in NR.}
  \label{fig:sm_hz_reduction}
\end{figure*}

\begin{figure*}[htbp]
  \centering
  \includegraphics[width=\textwidth]{figures/appendix_c/fig_hz4_violin.pdf}
  \caption{Run-level dispersion of $T_{0.95}^{(k)}$ and $\Sigma_s^{(k)}$ by
  hazard-source location. Points are individual runs, violins show
  distributional shape, and median bars summarise central tendency. The
  exposure distributions are more location-sensitive than the clearance-time
  distributions.}
  \label{fig:sm_hz_violin}
\end{figure*}

\begin{figure*}[htbp]
  \centering
  \includegraphics[width=\textwidth]{figures/appendix_c/fig_hz5_delta_bars.pdf}
  \caption{Improvement rates relative to Static routing for Density-60 and
  DH-60. Error bars propagate bootstrap confidence bounds. The risk-reduction
  advantage of DH-60 is consistent across locations and largest for the FR
  source class.}
  \label{fig:sm_hz_delta}
\end{figure*}

\subsection{Implications for Tail-Risk Dispersion}

The location analysis clarifies why tail-risk evaluation must report both
evacuation-time quantiles and cumulative-exposure indicators. Hazard-source
placement has a limited effect on $\mu_g[T_{0.95}]$ but a large effect on
$\mu_g[\Sigma_s]$ and $\mu_g[H]$. The FR class is the clearest example: it is not the
slowest class in tail clearance, yet it produces the highest Static exposure
and the largest incremental benefit from explicit hazard coupling. Therefore,
source distance to exits or main roads is not sufficient as a scalar
predictor of tail risk. What matters is the interaction between source
location, local path redundancy, crowd redistribution, and the time-varying
hazard field.

These results support the main paper's conclusion that density-only guidance
and density--hazard coupled guidance are not interchangeable. Density-60 is
an effective clearance-time intervention, but DH-60 is the risk-control
intervention: it preserves the clearance-time gains while suppressing
location-dependent exposure tails. In practical deployment, a detected
  hazard source remote from both exits and main roads may therefore warrant
  prioritised hazard-aware updating in future operational studies, because that
  is the source class with the largest incremental reduction in the model-based
  exposure indicators relative to density-only rerouting.



% =========================================================
% =========================================================
%
% Appendix D: Spatial Diagnostic Toolkit
%
% =========================================================
% =========================================================

\section{Spatial Diagnostic Toolkit}
\label{sm:spatial_diagnostics}

This appendix defines the spatial diagnostics used to localise persistent
congestion and hazard exposure in the main text
(Section~\ref{sec:evaluation}). Unlike system-level summary metrics, these
indicators reveal where evacuation risk accumulates and therefore support
bottleneck identification and interpretation.

\subsection{Route-use and density-load fields}

For each raster cell $\mathbf{p}$, route-use frequency is the number of
recorded agent positions falling in that cell:
\begin{equation}
F_{\mathrm{route}}(\mathbf{p})
= \sum_{n=1}^{N_t}\sum_{i\in\mathcal{A}_n}
\mathbb{I}\!\big(\mathbf{x}_i(t_n)\in \mathbf{p}\big).
\end{equation}

The cumulative density-load field integrates the raster density over time
without applying a congestion mask:
\begin{equation}
L_{\rho}(\mathbf{p})
= \sum_{n=1}^{N_t} \rho(\mathbf{p},t_n)\,\Delta t .
\end{equation}
This differs from $L_{\mathrm{cong}}(\mathbf{p})$, which accumulates
density only during slow-and-crowded states.

\subsection{Congestion indicators}

We define a joint congestion indicator on the raster:
\begin{equation}
\mathbb{I}_{\mathrm{cong}}(\mathbf{p}, t)
= \mathbb{I}\!\big(v(\mathbf{p}, t) \leq v_{\mathrm{th}}\big)
  \;\wedge\;
  \mathbb{I}\!\big(\rho(\mathbf{p}, t) \geq \rho_{\mathrm{th}}\big),
\end{equation}
where the reported analysis fixes the speed and density thresholds at
$v_{\mathrm{th}} = 0.5$\,m/s and
$\rho_{\mathrm{th}} = 3.0$\,ped/m$^2$, respectively.

\subsection{Aggregated congestion measures}

\begin{align}
T_{\mathrm{cong}}(\mathbf{p})
  &= \sum_{n=1}^{N_t}
     \mathbb{I}_{\mathrm{cong}}(\mathbf{p}, t_n)\, \Delta t,
  && \text{[s] \quad (congested duration)}, \\[4pt]
L_{\mathrm{cong}}(\mathbf{p})
  &= \sum_{n=1}^{N_t}
     \rho(\mathbf{p}, t_n)\,
     \mathbb{I}_{\mathrm{cong}}(\mathbf{p}, t_n)\, \Delta t,
  && \text{[(ped/m$^2$)$\cdot$s] \quad (congestion load)}.
\end{align}
$T_{\mathrm{cong}}$ measures persistence; $L_{\mathrm{cong}}$ measures
severity.

\subsection{Static-relative spatial summaries}

For each strategy $s$, the run-level fields are first averaged cellwise over
the $K=50$ matched scenarios,
\[
\overline{X}^{s}(\mathbf{p})
=\frac{1}{K}\sum_{k=1}^{K}X^{s,k}(\mathbf{p}).
\]
All subsequent summaries are computed from this ensemble-mean field; the
overbar is omitted below and in the main text for readability. The common
analysis domain is the walkable-road mask
\[
\mathcal{R}=\{\mathbf{p}:\mathrm{road}(\mathbf{p})>0\},
\]
obtained from the shared rasterised road layer (\texttt{layer\_2}). Using the
same fixed cell set for all strategies prevents off-network cells from
changing the Static-relative comparison.

For any non-negative spatial field $X(\mathbf{p})$, positive-cell
summaries are computed over
\[
\mathcal{P}^{+}_{X}=\{\mathbf{p}\in\mathcal{R}: X(\mathbf{p})>0\}.
\]
The positive-cell 95th percentile is denoted
$Q^{+}_{0.95}(X)$ and is computed from
$\{X(\mathbf{p}):\mathbf{p}\in\mathcal{P}^{+}_{X}\}$. For
Static-relative summaries, the same-field Static threshold is
\[
q^{\mathrm{Static}}_{0.99}(X)
= Q^{+}_{0.99}\!\left(X^{\mathrm{Static}}\right),
\]
and the number of cells above that threshold for strategy $s$ is
\[
N^{s}_{>0.99}(X)
=
\left|\{\mathbf{p}\in\mathcal{R}: X^{s}(\mathbf{p})
> q^{\mathrm{Static}}_{0.99}(X)\}\right|.
\]

\subsection{Bottleneck detection}

A cell is flagged as a candidate bottleneck if:
\begin{equation}
\mathbb{I}_{\mathrm{cand}}(\mathbf{p})
= \mathbb{I}\!\big(T_{\mathrm{cong}}(\mathbf{p}) \geq T_{\mathrm{th}}\big)
  \;\wedge\;
  \mathbb{I}\!\big(L_{\mathrm{cong}}(\mathbf{p}) \geq L_{\mathrm{th}}\big),
\end{equation}
where thresholds are set as high percentiles (e.g., 95th) of their
respective distributions. Candidate cells are clustered via
DBSCAN~\citep{ester1996density} into spatially distinct congestion zones.
This procedure extracts spatially coherent congestion zones that act as
bottlenecks in the evacuation process.

Together, these spatial diagnostics provide a fine-grained view of route use
and congestion, complementing the system-level performance metrics reported
in the main text. Soft Risk and Hard Harm remain the formal exposure outcomes
defined in \S\ref{sec:eval_exposure}.


\newpage
\section*{Acknowledgements}
\begin{sloppypar}
This work is supported by the Guangdong Provincial Project
(No.~2023QN10H717), the
Guangzhou--HKUST\allowbreak(GZ) Joint Funding Program
(No.~2025A03J3639), and the AI Research and Learning Base of Urban
Culture Project (No.~2023WZJD008).
\end{sloppypar}

\section*{Declaration of generative AI and AI-assisted technologies in the writing process} 
During the preparation of this work the authors used ChatGPT for language polishing. After using this tool, the authors reviewed and edited the content as needed and take full responsibility for the content of the publication.

% \section*{Declaration of Competing Interest}
% The authors declare that they have no known competing financial interests or
% personal relationships that could have appeared to influence the work reported
% in this paper.

% \section*{Data Availability}
% Upon publication, the derived Shipai raster environment, scenario parameter
% files, analysis scripts, and representative simulation outputs will be made
% available in a public repository. Source GIS or imagery layers subject to
% third-party restrictions will be represented by derived raster masks and
% metadata sufficient for reproducing the reported experiments.


\newpage
\bibliographystyle{elsarticle-harv}
\IfFileExists{reference.bib}{%
  \bibliography{reference}%
}{%
  \bibliography{paper/reference}%
}

\end{document}
